% ============================================================
% Chapter 8: China AI — The Eastern Front of the Intelligence Race
% Book: AI Figures: Who's Who in the Age of AI
% Language: Simplified Chinese (per series convention for China chapter)
% ============================================================

\chapter{China AI: The Eastern Front of the Intelligence Race}
\label{ch:china}

% China's AI ecosystem is not a copy of Silicon Valley — it is a parallel
% universe with its own logic, its own heroes, and its own rules.

如果说硅谷是全球 AI 革命的第一极，那么中国就是这场革命中最令人惊叹的
第二极——一个在芯片封锁、出口管制和地缘政治高压下依然持续产出
世界级成果的平行宇宙。

2025 年 1 月 20 日，DeepSeek-R1 的发布震动了整个科技世界。
一家来自杭州的量化对冲基金衍生公司，用不到 600 万美元的训练成本，
训练出了比肩 OpenAI o1 的推理模型。这一事件不仅引发了全球科技股的
剧烈震荡——NVIDIA 单日市值蒸发近 6000 亿美元——更从根本上颠覆了
"AI 竞赛 = 烧钱竞赛"的叙事。它证明了一个简单而深刻的事实：
\textbf{中国 AI 从来不只是"追赶者"，它正在成为规则的重新定义者}。

但 DeepSeek 只是冰山一角。在它背后，是一个由学术天才、连续创业者、
大厂实验室和国家战略共同编织的庞大生态系统。从清华姚班走出的年轻
创始人到微软亚研院的资深科学家，从 BAT 的万人 AI 团队到寒武纪兄弟的
芯片梦想，从 Andrew Ng 在百度种下的种子到汤晓鸥留下的遗产——
这一章讲述的是近百位塑造中国 AI 格局的关键人物的故事。

\begin{corebox}[本章结构]
本章按组织类型和技术方向组织 130+ 位中国 AI 关键人物：
\begin{itemize}[noitemsep]
  \item \textbf{DeepSeek}（§\ref{sec:deepseek}）：震惊世界的黑马
  \item \textbf{六小虎}（§\ref{sec:six-tigers}）：智谱、月之暗面、MiniMax、百川、零一万物、阶跃星辰
  \item \textbf{大厂 AI 实验室}（§\ref{sec:bigtech}）：字节、百度、阿里、腾讯、华为
  \item \textbf{AI 四小龙}（§\ref{sec:four-dragons}）：商汤、旷视、科大讯飞
  \item \textbf{学术领袖}（§\ref{sec:academics}）：从院士到海归教授
  \item \textbf{AI 芯片先驱}（§\ref{sec:chips}）：寒武纪、摩尔线程、壁仞科技
  \item \textbf{具身智能}（§\ref{sec:robotics}）：宇树、智元、银河通用、傅利叶
  \item \textbf{自动驾驶}（§\ref{sec:autonomous-driving}）：小马智行、文远知行、Momenta、蔚小理
  \item \textbf{开源生态与 AI 基础设施}（§\ref{sec:open-source}）：ColossalAI、硅基流动、面壁智能
\end{itemize}
\end{corebox}


% ============================================================
%  §8.1  DEEPSEEK
% ============================================================

\section{DeepSeek: The Dark Horse That Shook the World}
\label{sec:deepseek}

2025 年 1 月的"DeepSeek 时刻"是中国 AI 历史上的分水岭。
一家成立仅两年的公司，以极低的成本和极高的工程效率，
证明了在芯片受限条件下仍能训练出世界一流的大模型。
但要理解 DeepSeek 的成功，必须先理解它背后的人。

\subsection{梁文锋 (Liang Wenfeng) — 隐形的亿万富翁量化大师}

\begin{keybox}[人物档案：梁文锋]
\textbf{生年}：1985 年 \quad \textbf{籍贯}：广东 \quad \textbf{教育}：浙江大学电子信息工程 \\
\textbf{职务}：幻方量化联合创始人、DeepSeek 创始人兼 CEO \\
\textbf{荣誉}：2025 年 \textit{Time} AI 100、\textit{Fortune} 最具影响力商业人物
\end{keybox}

在 DeepSeek-R1 让全球科技圈震动之前，几乎没有人听说过梁文锋。
这正是他想要的。

梁文锋的故事始于量化金融。他在浙江大学攻读电子信息工程时，
就开始将机器学习方法应用于市场预测。毕业后，他与合伙人共同创立了
幻方量化（High-Flyer Quant），名字取自数学中的"幻方"（Magic Square）——
一个最早在中国被发现的数学对象。这个名字暗示了梁文锋的方法论核心：
\textbf{用数学的精确性驾驭不确定性}。

幻方量化管理着约 80 亿美元的资产，是中国顶级量化对冲基金之一。
2025 年，幻方的基金回报率高达 56.6\%，在中国大型量化基金中排名第二。
但真正让梁文锋与其他量化基金经理区分开的，是他在 2021 年做出的一个
在当时看来近乎疯狂的决定：\textbf{大规模购入 NVIDIA GPU}。

在美国对华芯片出口管制生效之前，梁文锋已经积累了数千块 A100 芯片。
这批硬件最初是为量化交易的 AI 模型准备的，但梁文锋的野心远不止于此。
2023 年，他创立了 DeepSeek，将这些算力资源投入到一个更宏大的目标：
\textbf{通用人工智能（AGI）}。

\begin{whybox}[为什么一个量化基金经理能做出世界级 AI？]
答案在于量化交易与 AI 研究之间被低估的协同效应：
\begin{enumerate}[noitemsep]
  \item \textbf{算力积累}：顶级量化基金本身就需要大规模 GPU 集群来训练
        交易模型，幻方在 2021 年之前就建立了 10000+ GPU 的算力基础。
  \item \textbf{工程文化}：量化交易要求极致的系统优化——延迟以微秒计，
        成本以 basis point 计。这种"榨干每一滴硬件性能"的文化，
        直接塑造了 DeepSeek "用 H800 做出 H100 效果"的技术路线。
  \item \textbf{人才磁场}：量化基金的高薪吸引了大量顶尖计算机科学毕业生，
        这些人可以无缝转向 AI 研究。
  \item \textbf{财务自由}：幻方的利润为 DeepSeek 提供了"不需要向 VC 交代"
        的财务独立性，使其可以做长期、高风险的基础研究。
\end{enumerate}
\end{whybox}

梁文锋的管理风格极为低调——他几乎从不接受媒体采访，不出席公开活动，
甚至没有公开的社交媒体账号。在 DeepSeek-R1 发布引发全球轰动后，
他依然保持沉默。这种极端的低调与 Sam Altman 的高调形成了鲜明对比，
也让外界对 DeepSeek 内部的运作方式充满好奇。

据 2025 年斯坦福胡佛研究所的一份报告分析，DeepSeek 的核心技术团队
约 140 人，平均年龄不到 30 岁，大多来自中国顶尖高校（北大、清华、
中科大、浙大）的应届或近期博士毕业生。这个团队规模仅为 OpenAI 的
二十分之一，却产出了比肩甚至超越前者的技术成果。

\subsection{DeepSeek 核心研究团队}

DeepSeek 的论文作者名单是了解其技术核心的窗口。与大多数 AI 实验室
不同，DeepSeek 的论文通常列出大量作者（V3 论文有 130+ 位作者），
但按字母序排列，不设"第一作者"或"通讯作者"——这反映了其
\textbf{扁平化、集体主义}的组织文化。

以下是从 DeepSeek-V3、R1 及 V3.2 论文中反复出现的关键研究者：

\paragraph{Aixin Liu（刘爱鑫）}
DeepSeek-V3 技术报告的名义第一作者（按字母序）。
在 V3、R1、V3.2 三篇核心论文中均出现，
是 DeepSeek 模型架构设计的核心人物之一。

\paragraph{Daya Guo（郭达雅）}
DeepSeek-R1 论文在 \textit{Nature} 正式发表时的第一作者，
Sun Yat-sen University（中山大学）校友。
同时也是 DeepSeekMath 和 DeepSeekCoder 系列的核心贡献者，
专注于代码智能和数学推理。他的 Google Scholar 主页显示其
研究横跨代码预训练模型（CodeBERT 的共同作者）和
数学推理（DeepSeekMath），是 DeepSeek 在推理能力方面的
技术骨干。

\paragraph{Chong Ruan（阮冲）}
DeepSeek 的联合创始人之一，在多篇核心论文中出现。

\paragraph{Damai Dai（戴大迈）}
DeepSeek MoE（Mixture of Experts）架构的核心设计者之一，
在 DeepSeekMoE 到 V3 的演进中扮演关键角色。

\paragraph{Bei Feng（冯贝）}
DeepSeek 基础设施和训练系统的核心工程师，
在 V3、R1、V3.2 中均有贡献，
负责分布式训练框架和通信优化（DualPipe、DeepEP 等）。

\begin{corebox}[DeepSeek 的"零 Loss Spike"工程奇迹]
DeepSeek-V3 的技术报告披露了一个令业界震惊的细节：
整个 671B 参数模型的预训练过程中实现了"zero irrecoverable loss spikes"——
即没有一次不可恢复的训练异常。考虑到 V3 使用的是
受出口管制限制的 H800 GPU（NVLink 带宽低于 H100），
这一成就需要极其成熟的容错基础设施和训练工程能力。

关键创新包括：
\begin{itemize}[noitemsep]
  \item \textbf{DualPipe}：通过更精细的计算-通信重叠减少 pipeline bubble
  \item \textbf{DeepEP}：为 MoE 的 expert parallelism 设计的专用 all-to-all 通信内核
  \item \textbf{FP8 混合精度训练}：在不损失精度的前提下大幅提升训练效率
\end{itemize}

这些工程创新使 DeepSeek 用 2048 块 H800 在约两个月内完成了 V3 的训练，
总成本约 557 万美元——不到 OpenAI 训练 GPT-4 估计成本的十分之一。
\end{corebox}

\subsection{DeepSeek 的技术演进路线}

理解 DeepSeek 的人才结构，需要了解其技术路线图：

\begin{enumerate}[noitemsep]
  \item \textbf{DeepSeekMoE}（2024 年初）：首次提出 Fine-Grained Expert
        Segmentation 和 Shared Expert Isolation，为后续 MoE 架构奠定基础。
  \item \textbf{DeepSeek-V2}（2024 年中）：引入 Multi-Head Latent Attention（MLA）
        和 DeepSeekMoE 架构，在推理效率上取得突破。
  \item \textbf{DeepSeek-V3}（2024 年 12 月）：671B 参数、37B 激活参数的
        MoE 模型，以极低成本达到 GPT-4 级别性能。
  \item \textbf{DeepSeek-R1}（2025 年 1 月）：通过纯强化学习（不需要
        人工标注的推理轨迹）激发 LLM 的推理能力，论文发表于 \textit{Nature}。
  \item \textbf{DeepSeek-V3.2}（2025 年 12 月）：引入 DeepSeek Sparse Attention（DSA），
        在 GPT-5 和 Gemini 3.0 Pro 级别上达到开放权重模型新高。
\end{enumerate}

每一步演进都代表着团队在特定技术方向上的突破。从 MoE 架构（V2/V3）
到强化学习推理（R1）再到注意力机制创新（V3.2），DeepSeek 展现出了
一个小而精团队在技术栈全方位的创新能力。


% ============================================================
%  §8.2  SIX LITTLE TIGERS
% ============================================================

\section{六小虎（Six Little Tigers）：中国大模型创业的第一梯队}
\label{sec:six-tigers}

"大模型六小虎"是中国 AI 创投圈对六家头部大模型创业公司的统称——
智谱 AI、月之暗面、MiniMax、百川智能、零一万物和阶跃星辰。
这六家公司均在 2023--2024 年间达到独角兽估值，
由中国最顶尖的 AI 技术人才创立，
背后站着阿里、腾讯、美团等科技巨头的战略投资。

截至 2026 年 3 月：智谱和 MiniMax 已在港交所上市，
月之暗面估值飙升至 180 亿美元，
阶跃星辰完成 50 亿元 B+ 轮融资，
而零一万物则转向了基于 DeepSeek 开源模型的应用层路线。

\begin{keybox}[六小虎一览（2026 年 3 月）]
\begin{tabular}{llllr}
\textbf{公司} & \textbf{创始人} & \textbf{成立} & \textbf{状态} & \textbf{估值/市值} \\
\hline
智谱 AI & 唐杰、张鹏 & 2019 & 港交所上市 & $\sim$1000 亿港元 \\
MiniMax & 闫俊杰 & 2021 & 港交所上市 & $\sim$1100 亿港元 \\
月之暗面 & 杨植麟 & 2023 & 私募 & $\sim$180 亿美元 \\
阶跃星辰 & 姜大昕 & 2023 & B+ 轮 & $\sim$350 亿元 \\
百川智能 & 王小川 & 2023 & A 轮 & $\sim$200 亿元 \\
零一万物 & 李开复 & 2023 & 转型应用层 & -- \\
\end{tabular}
\end{keybox}


% ---- 8.2.1 智谱 ----
\subsection{智谱 AI (Zhipu AI) — 清华血统的"大模型第一股"}

智谱 AI 是中国大模型创业公司中学术根基最深的一家。
它的故事始于 2006 年清华大学知识工程实验室的一个项目：AMiner——
一个学术社交网络分析与挖掘系统。近二十年后，这个项目的带头人
唐杰带领团队走出实验室，创建了中国第一家登陆港交所的
通用大模型公司。

\subsubsection{唐杰 (Tang Jie) — 躲在喧嚣之外的"真大佬"}

\begin{keybox}[人物档案：唐杰]
\textbf{生年}：$\sim$1979 年 \quad \textbf{教育}：清华大学计算机系博士 \\
\textbf{职务}：智谱 AI 首席科学家、清华大学计算机系教授 \\
\textbf{成就}：AMiner 创建者、GLM 系列模型奠基人 \\
\textbf{持股}：智谱 7.41\%，实际控制人之一
\end{keybox}

2006 年，清华计算机系的博士毕业生唐杰站在了一个人生岔路口：
出国还是留校。他选择了后者，效仿的对象是北大方正创始人王选——
以教授身份带领团队将实验室成果转化为产业奇迹。

唐杰在知识工程和社交网络挖掘领域深耕近二十年，
AMiner 系统累计服务了数百万科研工作者。
但真正让他走向产业前台的是 GPT-3 的出现。
2020 年，GPT-3 的发布让唐杰和团队确信大模型将是未来方向。
他们从 GLM（General Language Model）系列开始——
一种采用自回归填空预训练目标的架构，走出了与 GPT 不同的技术路线。

在智谱 2026 年 1 月 8 日上市当天，唐杰发布内部信，
宣布"全面回归基础模型研究"，并预告新一代 GLM-5 模型即将发布。
这封信的核心信息是：\textbf{上市不是终点，而是加速基础研究的起点}。

\subsubsection{张鹏 (Zhang Peng) — CEO 与商业化推手}

\begin{keybox}[人物档案：张鹏]
\textbf{教育}：清华大学计算机系学士、博士（知识图谱方向） \\
\textbf{职务}：智谱 AI CEO \\
\textbf{背景}：第 14 届北京市政协委员
\end{keybox}

张鹏是智谱的商业化引擎。他本科和博士均毕业于清华计算机系，
研究方向为知识图谱。与纯学术背景的唐杰互补，张鹏负责将
技术能力转化为可持续的商业模式。

智谱的商业模式是"双轮驱动"：高毛利的本地化部署（2025 上半年
贡献 85\% 收入，毛利率 59\%）为基石，云端 API 服务为增长引擎。
这一模式帮助智谱在 IPO 时获得了 1159 倍的公开认购超额。

\subsubsection{其他核心成员}

\paragraph{刘德兵 (Liu Debing)}
智谱董事长，中科院博士，师从高文院士。
负责战略及业务方向，与唐杰共同为实际控制人，
合计控制智谱 36.97\% 表决权。

\paragraph{李涓子 (Li Juanzi)}
清华大学计算机系教授，知识工程专家。
智谱创始团队成员，与唐杰同为清华知识工程实验室核心学者，
为智谱的知识图谱技术提供了学术根基。

\paragraph{许斌 (Xu Bin)}
清华大学教授，智谱创始团队成员，一致行动人之一。


% ---- 8.2.2 月之暗面 ----
\subsection{月之暗面 (Moonshot AI / Kimi) — AGI 追梦人}

月之暗面的名字取自 Pink Floyd 的专辑 \textit{The Dark Side of the Moon}——
"你通常只能看到月亮被照亮的一面，但暗面充满了未知与潜力。"
这正是其创始人杨植麟的 AGI 哲学。

\subsubsection{杨植麟 (Yang Zhilin) — 从 Transformer-XL 到 Kimi}

\begin{keybox}[人物档案：杨植麟]
\textbf{生年}：1992 年 \quad \textbf{籍贯}：广东汕头 \\
\textbf{教育}：清华大学计算机系学士、CMU 计算机科学博士 \\
\textbf{博士导师}：唐杰（清华）、Ruslan Salakhutdinov \& William Cohen（CMU） \\
\textbf{代表论文}：Transformer-XL、XLNet \\
\textbf{职务}：月之暗面创始人兼 CEO、清华大学助理教授
\end{keybox}

杨植麟的学术血统堪称完美。1992 年生于广东汕头，
本科毕业于清华大学计算机系，博士就读于 CMU——
师从深度学习大师 Ruslan Salakhutdinov 和 William Cohen。
在清华阶段，他的导师正是后来创立智谱的唐杰。

他的博士研究直接催生了两篇具有里程碑意义的论文：

\begin{itemize}[noitemsep]
  \item \textbf{Transformer-XL}（2019）：提出循环机制处理超长序列，
        解决了原始 Transformer 的固定上下文窗口限制。
        这一思想直接影响了后来所有长上下文模型的设计。
  \item \textbf{XLNet}（NeurIPS 2019 Oral）：结合自回归和自编码的优势，
        提出排列语言模型（Permutation Language Model），
        在发布时刷新了 18 项 NLP 任务的 SOTA。
\end{itemize}

这两篇论文的共同主题是\textbf{让模型处理更长的上下文}——
而这恰恰成为了 Kimi 的核心技术特征。

2023 年初，杨植麟在硅谷做了一个精确的计算：要创办一家
以 AGI 为目标的公司，需要在几个月内融资超过 1 亿美元。
这仅仅是"入场券"。十二个月后，这个数字膨胀了 13 倍。
他创立了月之暗面，以"长文本理解"作为技术切入点——
Kimi Chat 最初以支持 20 万字超长上下文而闻名，
后续的 Kimi K2 模型在 2025 年夏天跻身国际排行榜前十。

到 2026 年 3 月，月之暗面的估值已从创立时的不到 10 亿美元
飙升至 180 亿美元——不到 90 天内完成三轮融资、合计超 22 亿美元。

\begin{whybox}[Transformer-XL 到 Kimi 的技术连续性]
杨植麟创立月之暗面并非偶然的商业决策，而是其学术研究的自然延伸。
Transformer-XL 解决了 Transformer 的"上下文遗忘"问题，
XLNet 将这一思想推向了更强大的预训练框架。
Kimi 的长上下文能力正是这条技术线索的产品化——
从论文到产品，用了整整五年时间。
这是中国 AI 创业中"学术驱动型创业"的典型案例。
\end{whybox}

\subsubsection{Zihang Dai（戴自航）— Transformer-XL 的另一位共同作者}

Zihang Dai 是 Transformer-XL 和 XLNet 的共同第一作者，
同样是 CMU 博士（导师 Yiming Yang）。
博士毕业后，他加入 Google Brain 担任研究科学家，
参与了 Gemini 等模型的开发。
2024 年，Dai 成为 xAI（Elon Musk 的 AI 公司）的联合创始人之一。
2026 年 3 月，据报道 Dai 已从 xAI 离职，
成为该公司持续的联合创始人离职潮中的一员。
虽然 Dai 没有直接加入月之暗面，
但他与杨植麟的学术合作关系代表了中国 AI 人才在
全球顶级实验室之间流动的典型路径。

\subsubsection{Guokun Lai（赖国坤）}

CMU 博士（导师 Yiming Yang），清华本科。
2024 年加入月之暗面担任研究工程师。
作为 Transformer-XL 的合著者之一，他的加入标志着
杨植麟正在将博士期间的核心合作网络引入公司。


% ---- 8.2.3 MiniMax ----
\subsection{MiniMax — 从商汤实习生到港股 IPO 纪录}

MiniMax（上海稀宇科技）是六小虎中全球化程度最高的一家。
其旗舰产品海螺 AI/Talkie 覆盖全球 200 多个国家、
触达 2.12 亿用户，海外收入占比超过 70\%。
2026 年 1 月 9 日，MiniMax 在港交所上市，
公开认购超额 1837 倍，上市首日暴涨超 100\%，
市值一度突破 3800 亿港元。

\subsubsection{闫俊杰 (Yan Junjie) — 河南小伙的 AI 帝国}

\begin{keybox}[人物档案：闫俊杰]
\textbf{生年}：$\sim$1990 年 \quad \textbf{籍贯}：河南 \\
\textbf{教育}：中国科学技术大学博士（26 岁获得） \\
\textbf{前职}：商汤科技高级研究员 \\
\textbf{职务}：MiniMax 创始人兼 CEO \\
\textbf{财富}：$\sim$258 亿元（2026 年 1 月）
\end{keybox}

闫俊杰的故事是中国 AI 创业中最具戏剧性的之一。
出生于河南一个普通家庭，闫俊杰从小展现出过人的数理天赋。
他在中国科学技术大学一路攻读至博士，26 岁即获得博士学位——
研究方向集中在计算机视觉和深度学习。
博士毕业后，闫俊杰加入商汤科技，师从汤晓鸥团队体系，
在商汤担任高级研究员期间，他不仅在技术层面积累了深厚功底，
更展现出了卓越的人才培养能力——他指导的多位实习生
后来成为华为"天才少年"计划的入选者。
在商汤的经历让闫俊杰深刻认识到两点：
(1) 大模型将比计算机视觉有更广阔的应用空间；
(2) 中国需要自己的通用大模型公司。
2021 年，32 岁的闫俊杰创立 MiniMax（上海稀宇科技），
正式开启大模型创业。他的创业时间节点值得注意——
比 ChatGPT 发布早了整整一年半，比六小虎中的其他公司
都更早下注大模型方向。这一"先发优势"使 MiniMax 在
产品打磨和全球化布局上赢得了宝贵的时间窗口。

MiniMax 的差异化在于两点：
\begin{enumerate}[noitemsep]
  \item \textbf{B+C 双轮驱动}：既有面向开发者的 API 业务，
        又有直接面向消费者的 AI 伴侣产品（Talkie/海螺 AI）。
  \item \textbf{全球化优先}：与其他主要面向中国市场的六小虎不同，
        MiniMax 从一开始就将海外市场作为核心战略方向，
        Talkie 在全球社交 AI 领域占据了显著市场份额。
\end{enumerate}

闫俊杰在 2026 年世界人工智能大会上的发言揭示了他对行业的判断：
"在 AI 时代，企业成功的终极决定因素不是资金和资源的投入，
而是智能能力提升的速度——这取决于研发效率。"

\begin{warnbox}[MiniMax 的盈利困境]
尽管市值一度超过千亿港元，MiniMax 仍处于大幅亏损状态。
2025 年前 9 个月亏损超过 35 亿元，约为收入的近 10 倍。
近三年累计亏损预计超过 90 亿元。
创始人财富与公司盈利之间的巨大鸿沟，
是当前中国大模型公司共同面临的挑战。
\end{warnbox}


% ---- 8.2.4 百川 ----
\subsection{百川智能 (Baichuan AI) — 搜狗老将的"二次创业"}

\subsubsection{王小川 (Wang Xiaochuan) — 从奥赛金牌到 AI 创业}

\begin{keybox}[人物档案：王小川]
\textbf{生年}：1978 年 10 月 3 日 \quad \textbf{籍贯}：四川成都 \\
\textbf{教育}：清华大学计算机系学士、硕士、工程博士、EMBA \\
\textbf{荣誉}：1996 年 IOI（国际信息学奥赛）金牌 \\
\textbf{前职}：搜狗创始人兼 CEO（2010--2021） \\
\textbf{职务}：百川智能创始人兼 CEO
\end{keybox}

王小川是中国互联网史上最具传奇色彩的技术型创始人之一。
7 岁开始写程序，12 岁在中学生编程竞赛中获奖，
1996 年代表中国参加国际信息学奥林匹克竞赛（IOI）夺得金牌。
进入清华计算机系后，他一路从学士读到工程博士。

在搜狗，王小川开创了"三级火箭模式"——
用搜狗输入法 $\to$ 搜狗浏览器 $\to$ 搜狗搜索的产品链条，
成为追赶者破局搜索市场的唯一成功案例。
2017 年，他带领搜狗在纽交所上市。2021 年搜狗被腾讯收购后，
王小川短暂离开公众视野。

2023 年 4 月 10 日，王小川宣布创立百川智能。
凭借"IOI 金牌得主 + 清华 + 搜狗 CEO"的三重光环，
百川智能创下了中国大模型初创企业最快达到独角兽估值的纪录——
成立仅半年即估值超过 10 亿美元。

百川智能的差异化在于\textbf{AI 健康方向}。
王小川在多次公开演讲中强调，百川的长期愿景是成为
"AI 时代的健康顾问"——利用大模型理解医学知识、
辅助健康管理和医疗决策。2024 年 5 月推出的"百小应"
AI 助手，以及 2024 年 7 月 WAIC 上展示的"AI 健康顾问"，
都体现了这一战略方向。

2024 年 7 月，百川智能完成 A 轮融资，总金额约 50 亿元，
并计划以 200 亿元估值开启 B 轮。


% ---- 8.2.5 零一万物 ----
\subsection{零一万物 (01.AI) — 李开复的"最后一次创业"}

\subsubsection{李开复 (Kai-Fu Lee) — AI 传教士的创业之路}

\begin{keybox}[人物档案：李开复]
\textbf{生年}：1961 年 12 月 3 日 \quad \textbf{籍贯}：台湾 \\
\textbf{教育}：哥伦比亚大学计算机科学学士、CMU 计算机科学博士 \\
\textbf{前职}：Apple（语音识别）$\to$ SGI 副总裁 $\to$ Microsoft 副总裁 $\to$ Google 中国区总裁（2005--2009） \\
\textbf{著作}：\textit{AI Superpowers}、\textit{AI 2041} \\
\textbf{职务}：零一万物创始人兼 CEO、创新工场董事长
\end{keybox}

李开复是中国 AI 领域最知名的国际面孔。
他在 CMU 的博士研究（语音识别）是 AI 早期的重要工作，
此后在 Apple、SGI、Microsoft、Google 四家科技巨头担任高管，
每一段经历都留下了深刻印记——
在 Google，他一手建立了 Google 中国研究院，
培养了一批后来成为中国 AI 产业中坚力量的人才。

2009 年离开 Google 后，李开复创立了创新工场（Sinovation Ventures），
投资孵化了数百家科技初创公司。
他的畅销书 \textit{AI Superpowers}（2018）是第一本系统比较
中美 AI 竞争格局的英文著作，对西方世界理解中国 AI 产业
产生了深远影响。

2023 年，62 岁的李开复创立零一万物（01.AI），
被外界解读为他的"最后一次创业"。零一万物最初走的是
自研大模型路线——Yi 系列模型发布后获得了不错的社区反响。
然而，2025 年 DeepSeek 的开源攻势从根本上改变了李开复的策略。
在一次 SCMP 采访中，他坦言：
"DeepSeek 的出现让中美 AI 差距缩小到了三个月以内。"

此后，零一万物开始\textbf{利用 DeepSeek 的开源模型为企业客户
提供 AI 解决方案}——从自研基座模型转向应用层和企业服务。
这一转型虽然务实，但也标志着六小虎中第一家主动放弃
通用大模型预训练赛道的公司。


% ---- 8.2.6 阶跃星辰 ----
\subsection{阶跃星辰 (StepFun) — 微软老将的新征程}

\subsubsection{姜大昕 (Jiang Daxin) — 从 Bing 到 AGI}

\begin{keybox}[人物档案：姜大昕]
\textbf{教育}：纽约州立大学布法罗分校计算机科学博士 \\
\textbf{前职}：微软全球副总裁（2023 年 3 月升任），\\
\quad 微软亚洲互联网工程院（STCA）副院长兼首席科学家 \\
\textbf{职务}：阶跃星辰联合创始人兼 CEO \\
\textbf{代表工作}：CodeBERT、GraphCodeBERT、WizardLM、Bing 智能问答
\end{keybox}

姜大昕在微软的 16 年职业生涯横跨研究与工程。
2007 年加入微软亚洲研究院（MSRA）担任首席研究员，
2011 年转入 STCA（微软亚洲互联网工程院），
2017 年成为微软全球合伙人，
2023 年 3 月升任微软副总裁——离职前管理着 400 多人的团队，
负责 Bing 全球搜索体验和 NLP 技术研发。

他的学术贡献同样突出：CodeBERT 和 GraphCodeBERT
（代码预训练模型的奠基性工作）、
WizardLM（复杂指令跟随模型）的共同作者。
2008 年获得 SIGKDD 最佳应用论文奖。
2019 年，他带领团队将 Bing 智能问答功能从十几种语言
扩展至全球 100 多种语言和 200 个地区。

2023 年 4 月，姜大昕创立阶跃星辰，
与朱亦博（前 MSRA 研究员）和焦冰心共同创业。
2026 年 1 月，阶跃星辰完成超 50 亿元 B+ 轮融资，
创下过去 12 个月中国大模型赛道单笔最高融资纪录。
投资方包括多家国资机构（上国投、国寿股权、浦东创投等），
以及腾讯、启明、五源等老股东。

\subsubsection{张祥雨 (Zhang Xiangyu) — 从旷视到阶跃}

张祥雨曾任旷视科技研究院院长，是 AI 视觉领域的资深专家。
加入阶跃星辰后担任核心技术负责人，
为公司在多模态模型方向提供了关键技术能力——
阶跃星辰已发布 22 款自研基座模型，其中 16 款为多模态模型。

\subsubsection{印奇 (Yin Qi) — 意外的董事长}

2026 年 1 月，在阶跃星辰 B+ 轮融资的同时，
旷视科技联合创始人\textbf{印奇出任阶跃星辰董事长}——
这一人事任命在行业内引发广泛关注。
一位 AI 四小龙的创始人"空降"到六小虎，
既反映了中国 AI 产业从计算机视觉时代向大模型时代的
人才迁移，也暗示了两家公司在技术和资本层面的深度绑定。


% ============================================================
%  §8.3  BIG TECH AI LABS
% ============================================================

\section{大厂 AI 实验室：BAT、字节与华为}
\label{sec:bigtech}

中国大厂的 AI 实验室是另一个平行宇宙。
它们拥有六小虎无法匹敌的数据、分发渠道和资本，
但也面临着大企业特有的组织惯性和人才流失。


% ---- 8.3.1 字节跳动 ----
\subsection{字节跳动 AI / 豆包 (ByteDance / Doubao)}

字节跳动是中国 AI 应用层最成功的玩家。
其 AI 聊天产品"豆包"在 2025 年 12 月突破 1 亿 DAU，
成为全球日活最高的 AI 聊天应用之一——
而且据内部人士透露，豆包的用户增长和营销投入
是字节所有破亿 DAU 产品中最低的。

\subsubsection{张一鸣 (Zhang Yiming) — AI 基因的播种者}

\begin{keybox}[人物档案：张一鸣]
\textbf{生年}：1983 年 \quad \textbf{籍贯}：福建龙岩 \\
\textbf{教育}：南开大学软件工程学士 \\
\textbf{职务}：字节跳动创始人（2021 年卸任 CEO，半退休状态）
\end{keybox}

张一鸣虽然已于 2021 年卸任字节跳动 CEO，
但他为字节注入的"算法驱动一切"的基因是公司 AI 能力的根源。
今日头条是全球最早大规模使用推荐算法的内容平台之一，
TikTok 的推荐引擎更是重新定义了短视频产业。

字节跳动早在 2016 年就成立了 AI Lab，
并在计算机视觉、自然语言处理、语音处理等方向进行了深入布局。
2023 年，字节成立了专门的 \textbf{Seed 团队}（种子团队），
负责基础大模型的研发——这就是豆包背后的技术引擎。

2026 年 2 月，字节在美国开放了近 100 个 AI 岗位，
主要面向 Seed 团队，显示出其正在全球范围内
与 OpenAI、Google 等竞争顶尖 AI 人才。

\subsubsection{李航 (Li Hang) — 字节 AI Lab 的奠基者}

李航是字节跳动 AI Lab 的早期核心人物，
曾任微软亚洲研究院常务副院长，
是自然语言处理领域的资深专家。
他在字节期间主导了早期的 NLP 技术体系建设，
为后来的大模型能力奠定了基础。

\subsubsection{Seed 团队与豆包}

豆包的技术底座是字节自研的云雀大模型家族，
采用混合专家模型（MoE）架构。
Seed 团队在 2025--2026 年间发布了一系列模型：
Seed 1.8、Seed 2.0（多模态理解全面升级）、
Seedance 2.0（视频生成）、Seedream 5.0 Lite（图像生成）、
Seed LiveInterpret 2.0（实时翻译）。

这一产品矩阵的覆盖范围之广，在中国大厂 AI 实验室中首屈一指。


% ---- 8.3.2 百度 ----
\subsection{百度 AI / 文心一言 (Baidu / ERNIE)}

百度是中国最早全面押注 AI 的互联网巨头。
李彦宏常说："百度是一家 AI 公司，搜索只是 AI 的应用之一。"
这句话在 2015 年说出时被视为营销口号，
但十年后回头看，百度确实在 AI 上投入了超过 1000 亿元研发费用。

\subsubsection{李彦宏 (Robin Li) — 中国 AI 最早的信徒}

\begin{keybox}[人物档案：李彦宏]
\textbf{生年}：1968 年 11 月 17 日 \quad \textbf{籍贯}：山西阳泉 \\
\textbf{教育}：北京大学信息管理学士、University at Buffalo 计算机硕士 \\
\textbf{前职}：Infoseek 工程师、IDD Information Services 顾问 \\
\textbf{职务}：百度联合创始人、董事长兼 CEO \\
\textbf{专利}：超链分析（Hyperlink Analysis）——与 Google PageRank 独立提出
\end{keybox}

李彦宏在 2012 年就开始布局深度学习——
那一年 AlexNet 刚在 ImageNet 上引发深度学习革命。
他做出的三个关键决策塑造了百度的 AI 轨迹：

\begin{enumerate}[noitemsep]
  \item \textbf{2013 年}：成立百度深度学习研究院（IDL），
        这是中国互联网公司中最早的深度学习专门机构。
  \item \textbf{2014 年}：重金聘请 Andrew Ng（吴恩达）担任
        百度首席科学家——这一举动被视为中国 AI 产业国际化的标志性事件。
  \item \textbf{2019 年}：发布 ERNIE（文心大模型的前身），
        走"知识增强"的技术路线。
\end{enumerate}

然而，百度的 AI 战略在商业化层面长期面临"叫好不叫座"的困境。
文心一言（ERNIE Bot）用户规模虽然达到 4.3 亿，
日均调用量超过 15 亿次，但始终未能孵化出现象级 C 端应用，
B 端落地进展也较为缓慢。

2025 年 11 月，李彦宏做出了一个重大组织调整：
新设基础模型研发部与应用模型研发部，直接向 CEO 汇报——
绕过了 CTO 王海峰。这一调整被解读为
\textbf{李彦宏对"技术有积累、市场无响动"的不满}。

\subsubsection{王海峰 (Wang Haifeng) — 学术与产业的错位}

\begin{keybox}[人物档案：王海峰]
\textbf{职务}：百度 CTO、AI Group 及百度研究院负责人 \\
\textbf{荣誉}：ACL 历史上首位华人主席，三次入围工程院院士增选 \\
\textbf{成就}：主导百度知识图谱、飞桨平台、百度大脑、文心大模型
\end{keybox}

王海峰拥有中国 AI 产业界最亮眼的学术资历之一。
作为 ACL 首位华人主席，他主导构建了百度从芯片（昆仑芯）
到框架（飞桨 PaddlePaddle）到模型（文心）的全栈 AI 技术体系，
牵头获得国家技术发明奖、科技进步奖等多项荣誉。

然而，"纸面成绩"与市场表现之间的鸿沟成为他的阿喀琉斯之踵。
飞桨虽号称"中国首个产业级深度学习平台"，
却未能在开发者市场上与 PyTorch 竞争。
文心大模型已迭代至 5.0 版本，
但在用户心智中的存在感远逊于 DeepSeek 和 Kimi。

\subsubsection{Andrew Ng（吴恩达）— 播下的种子}

\begin{keybox}[人物档案：Andrew Ng]
\textbf{生年}：1976 年 \quad \textbf{教育}：CMU 本科、MIT 硕士、UC Berkeley 博士 \\
\textbf{前职}：Stanford 教授 $\to$ Google Brain 联合创始人 $\to$ \textbf{百度首席科学家（2014--2017）} \\
\textbf{现职}：Stanford 兼职教授、Landing AI CEO、DeepLearning.AI 创始人、Coursera 联合创始人
\end{keybox}

Andrew Ng 在百度的三年（2014--2017）虽然短暂，
但其影响深远。他在硅谷建立了百度 AI 研究院的美国分部，
引入了国际顶尖的研究方法论和人才标准。
更重要的是，他的到来向全球传递了一个信号：
\textbf{中国互联网公司正在认真对待 AI 研究}。

2017 年 3 月，Ng 在 Medium 上发表告别信，
表示"百度的 AI 实力令人难以置信，团队从上到下都充满人才"。
离开百度后，他创立了 Landing AI（专注制造业 AI）
和 DeepLearning.AI（AI 教育平台），
同时继续在 Stanford 和 Coursera 上推广 AI 教育——
他在 Coursera 上的机器学习课程已有超过 500 万学生注册。


% ---- 8.3.3 阿里巴巴 ----
\subsection{阿里巴巴 AI / 通义千问 (Alibaba / Qwen)}

阿里巴巴的 AI 布局以达摩院（DAMO Academy）为核心，
通义千问（Qwen）系列模型是其旗舰产品。
Qwen 团队是中国开源大模型生态中最活跃的贡献者之一——
开源模型超 400 个，衍生模型突破 20 万个，下载量超 10 亿次。

\subsubsection{周靖人 (Zhou Jingren) — 从微软到通义}

\begin{keybox}[人物档案：周靖人]
\textbf{前职}：Microsoft（长期任职） \\
\textbf{职务}：阿里云 CTO、通义实验室（Tongyi Lab）负责人 \\
\textbf{角色}：Qwen 大模型的最高技术负责人
\end{keybox}

周靖人从微软加入阿里巴巴后，一手搭建了通义大模型的技术体系。
2026 年 3 月，在 Qwen 技术负责人林俊阳（Lin Junyang）突然离职后，
周靖人直接接管了 Qwen 模型的顶层管理——
显示出阿里对这一核心 AI 资产的高度重视。

\subsubsection{林俊阳 (Lin Junyang) — Qwen 技术负责人的突然离职}

2026 年 3 月 4 日凌晨，就在 Qwen 3.5 小模型系列发布仅 24 小时后，
通义千问大模型技术负责人林俊阳在社交媒体上发布了一条简短消息，
随后被确认已于 3 月 3 日下午正式向阿里提交辞呈。

这一消息在 Qwen 团队内部引发了强烈情绪反应。
据 VentureBeat 报道，林俊阳的离职可能不是孤例——
"Qwen 团队中的关键人物正在离开"。
这一事件凸显了中国大厂 AI 实验室面临的核心挑战：
\textbf{如何在大企业体制内留住顶尖 AI 人才}。

\begin{whybox}[大厂 AI 实验室的人才困境]
中国大厂 AI 实验室面临一个结构性矛盾：
它们能提供六小虎无法匹敌的数据、算力和分发渠道，
但在激励机制上却难以与创业公司竞争。
一位在六小虎工作的核心研究者，可以通过股权激励
在公司上市后获得数亿乃至数十亿的财富回报。
而在大厂 AI 实验室，即使做出了世界级成果，
回报上限也被大企业的薪酬体系所约束。
林俊阳的离职、阿里的紧急管理层调整，
以及此前多家大厂 AI 团队的人才流失，
都是这一结构性矛盾的直接体现。
\end{whybox}


\begin{warnbox}[大厂系 AI 创业的悲剧样本：光年之外]
2023 年 2 月，美团联合创始人\textbf{王慧文 (Wang Huiwen)} 高调宣布创立
"光年之外"（Light Years Beyond），自封"中国的 OpenAI"。
凭借他在美团的传奇创业经历和顶级人脉，
光年之外在四个月内获得红杉中国（HongShan）等机构 2800 万美元投资，
美团 CEO 王兴个人也出资 500 万美元。

然而，仅四个月后，故事急转直下。
2023 年 6 月，王慧文因健康原因（据报道为抑郁症）
辞去所有公司职务。美团随即以 2.337 亿美元现金收购光年之外，
同时承接其 5066 万美元债务。
2024 年 11 月，据 36Kr 报道，王慧文健康恢复后重返美团，
领导一个独立的 AI 探索团队（GN06），
开发 AI 伴侣产品"Wow"。

光年之外的故事是中国 AI 创业热潮中最令人唏嘘的插曲——
它揭示了即便拥有顶级资源和人脉，
\textbf{创始人的身心健康仍是创业公司最脆弱的单点故障}。
\end{warnbox}


% ---- 8.3.4 腾讯 ----
\subsection{腾讯 AI / 混元 (Tencent / Hunyuan)}

腾讯在 AI 大模型竞赛中的策略是"低调投资 + 自研并行"。
一方面，腾讯投资了六小虎中的五家（除零一万物外）；
另一方面，腾讯自研了混元（Hunyuan）大模型和元宝 AI 应用。
2026 年春节期间，腾讯元宝与字节豆包、阿里千问三款应用
的营销战总花费超百亿元，共同跻身"月活亿级俱乐部"。

\subsubsection{张正友 (Zhang Zhengyou) — "张氏标定法"的发明者}

\begin{keybox}[人物档案：张正友]
\textbf{生年}：$\sim$1965 年 \quad \textbf{教育}：浙江大学电子工程学士、\\
\quad 法国南锡大学计算机硕士、巴黎第十一大学博士（1990）、理学博士（1994） \\
\textbf{前职}：Microsoft Research Redmond（20 年） \\
\textbf{职务}：腾讯首席科学家、腾讯 AI Lab 及 Robotics X 实验室主任 \\
\textbf{荣誉}：ACM Fellow、IEEE Fellow、2013 年 ICCV Helmholtz 时间考验奖
\end{keybox}

张正友是世界级的计算机视觉专家。
他在 1998--2000 年间提出的摄像机标定方法
（"张氏标定法"，Zhang's Camera Calibration）
是计算机视觉领域被引用最多的论文之一，
至今仍被全球广泛采用。2013 年，这一工作获得了
IEEE ICCV 颁发的 Helmholtz 时间考验奖。

在微软研究院工作 20 年后，张正友于 2018 年加入腾讯，
被聘为腾讯历史最高专业职级——17 级研究员/杰出科学家。
他在腾讯的角色横跨 AI 研究和机器人——
主导 AI Lab 的基础研究（包括混元大模型的技术方向），
同时领导 Robotics X 实验室探索下一代人机协作。


% ---- 8.3.5 华为 ----
\subsection{华为 AI / 昇腾与盘古 (Huawei / Ascend \& PanGu)}

华为在 AI 领域的定位与其他大厂截然不同：
它不追求消费级 AI 应用，而是专注于 \textbf{AI 基础设施}——
芯片（昇腾 Ascend）、计算框架（MindSpore）和行业大模型（盘古 PanGu）。
在美国芯片出口管制的压力下，华为的昇腾芯片成为中国 AI 算力
"去 NVIDIA 化"的核心选项。

\subsubsection{徐直军 (Xu Zhijun) — 华为 AI 战略的掌舵者}

\begin{keybox}[人物档案：徐直军]
\textbf{生年}：1967 年 \quad \textbf{籍贯}：安徽 \\
\textbf{教育}：东南大学通信工程学士 \\
\textbf{职务}：华为副董事长、轮值董事长 \\
\textbf{角色}：华为 AI 战略和昇腾芯片路线图的最高负责人
\end{keybox}

徐直军是华为 AI 战略的对外发言人和决策者。
在 2025 年华为全联接大会（HC 2025）上，他做了一场罕见的
50 分钟主题演讲，首次完整公布了昇腾芯片未来三年的迭代路线：
\begin{itemize}[noitemsep]
  \item \textbf{2026 年 Q1}：Ascend 950 系列（低精度数据格式突破、
        向量算力提升、互联带宽增加、自研 HBM）
  \item \textbf{2027 年}：Ascend 960
  \item \textbf{2028 年}：Ascend 970
\end{itemize}
"一年一代、算力翻倍"的路线图直指 NVIDIA——
徐直军明确表示："我们对为人工智能的长期快速发展
提供可持续且充裕算力充满信心。"

华为在昇腾生态上做出了四项关键决定：
坚持硬件变现、CANN 编译器开放、Mind 工具链全面开源、
openPangu 基础大模型全面开源。
这些举措的核心逻辑是：
\textbf{用开源生态换市场份额，用市场份额养芯片迭代}。

\subsubsection{张平安 (Zhang Ping'an) — 华为云与盘古大模型}

华为常务董事、华为云计算 CEO 张平安负责将昇腾算力和
盘古大模型转化为云服务产品。2025 年 6 月，他宣布
基于 CloudMatrix 384 超节点的新一代昇腾 AI 云服务全面上线，
同时发布盘古大模型 5.5——覆盖 NLP、计算机视觉、多模态、
预测和科学计算五大基础模型。

\begin{corebox}[华为的"超节点 + 集群"战略]
华为的 AI 算力战略可以浓缩为五个字："超节点 + 集群"。

由于受美国出口管制限制，华为无法获取最先进的制造工艺，
单颗芯片的性能上限被压制。华为的应对策略是通过
\textbf{系统级创新弥补单芯片劣势}：将 384 颗昇腾 NPU 和
192 颗鲲鹏 CPU 通过全新高速网络（灵衢互联协议）全对等互联，
形成一台"超级 AI 服务器"——单卡推理吞吐量达到
2300 tokens/s，相比非超节点提升近 4 倍。

这一思路与 DeepSeek 用软件创新弥补 H800 硬件劣势的逻辑一脉相承，
代表了中国 AI 在芯片受限条件下的共同策略：
\textbf{当单点不如对手时，用系统级设计赢回差距}。
\end{corebox}


% ============================================================
%  §8.4  FOUR DRAGONS & LEGACY AI
% ============================================================

\section{AI 四小龙与行业先驱}
\label{sec:four-dragons}

在大模型时代到来之前，中国 AI 产业的代名词是"AI 四小龙"——
商汤科技、旷视科技、依图科技和云从科技。
这四家公司以计算机视觉为核心技术，在安防、金融、出行等
垂直领域建立了庞大的商业版图。
大模型浪潮对它们既是挑战也是机遇。


% ---- 8.4.1 商汤 ----
\subsection{商汤科技 (SenseTime) — 失去灵魂人物之后}

\subsubsection{汤晓鸥 (Tang Xiaoou, 1968--2023) — 中国 CV 的教父}

\begin{keybox}[人物档案：汤晓鸥]
\textbf{生卒}：1968--2023 年 12 月 15 日 \\
\textbf{教育}：中国科学技术大学学士、Rochester 大学硕士、MIT 博士 \\
\textbf{职务}：商汤科技创始人兼执行董事、CUHK 信息工程系教授、上海人工智能实验室主任 \\
\textbf{荣誉}：IEEE Fellow，中国计算机视觉领域最具影响力的学者之一
\end{keybox}

2023 年 12 月 15 日深夜，汤晓鸥因病去世，享年 55 岁。
他的去世在中国 AI 界引发了巨大的震动和悲痛。

汤晓鸥的一生横跨学术与产业。在 CUHK（香港中文大学），
他建立了多媒体实验室（mmlab），培养了一批后来塑造了
中国计算机视觉产业的学生——包括商汤 CEO 徐立、
CTO 王晓刚，以及后来创立思谋科技的贾佳亚。
他的实验室是中国 AI 界的"黄埔军校"。

2014 年，汤晓鸥团队在人脸识别领域首次超越人类水平——
这一成就直接催生了商汤科技的创立。
商汤后来成长为全球估值最高的 AI 公司之一（最高时超过 120 亿美元），
并于 2021 年在港交所上市。

汤晓鸥去世后，商汤的股价大幅下跌——
市场对这家失去灵魂人物的公司的前景感到担忧。

\begin{whybox}[汤晓鸥的遗产：CUHK mmlab 的人才树]
汤晓鸥在 CUHK 建立的多媒体实验室（mmlab）
是中国 AI 人才培养最成功的案例之一。
从这个实验室走出的学生和博士后包括：
\begin{itemize}[noitemsep]
  \item \textbf{徐立}（Xu Li）——商汤科技 CEO
  \item \textbf{王晓刚}（Wang Xiaogang）——商汤科技 CTO
  \item \textbf{贾佳亚}（Jia Jiaya）——CUHK/HKUST 教授，思谋科技创始人
  \item \textbf{何恺明}（Kaiming He）——ResNet 发明者（虽非直接学生，但同为 CUHK 校友）
\end{itemize}
mmlab 的开源贡献（OpenMMLab 系列工具包）
至今仍是全球计算机视觉社区最广泛使用的基础设施之一。
\end{whybox}

\subsubsection{徐立 (Xu Li) — 接棒者}

\begin{keybox}[人物档案：徐立]
\textbf{教育}：CUHK 博士（师从汤晓鸥） \\
\textbf{职务}：商汤科技创始人、董事会主席兼 CEO \\
\textbf{角色}：汤晓鸥去世后，全面接管公司运营和战略
\end{keybox}

徐立是汤晓鸥在 CUHK 培养的核心学生之一。
在汤晓鸥去世后，他全面接管了商汤的运营，
带领公司从传统的计算机视觉赛道向大模型时代转型——
商汤的日日新（SenseNova）大模型和 SenseCore AI 基础设施
是这一转型的核心产品。

\subsubsection{王晓刚 (Wang Xiaogang) — 学术桥梁}

王晓刚同样出自汤晓鸥门下，担任商汤 CTO。
他在 CUHK 和商汤之间架起了学术与产业的桥梁，
主导了商汤在基础视觉研究方面的技术积累。


% ---- 8.4.2 旷视 ----
\subsection{旷视科技 (Megvii / Face++) — 清华姚班三剑客}

\subsubsection{印奇 (Yin Qi)、唐文斌 (Tang Wenbin)、杨沐 (Yang Mu)}

\begin{keybox}[人物档案：旷视三创始人]
\textbf{印奇}：1988 年生，清华姚班学士、哥伦比亚大学硕士。CEO。\\
\textbf{唐文斌}：1987 年生，清华姚班学士、硕士（知识工程组）。CTO。\\
\textbf{杨沐}：1988 年生，清华姚班学士、硕士。\\
三人均有 Microsoft Research Asia 或 Google China 实习经历。
\end{keybox}

旷视的故事始于清华大学的"姚班"——
由图灵奖得主姚期智创立的计算机科学实验班。
印奇、唐文斌、杨沐三人在姚班结识，
2011 年共同创立旷视科技，以人脸识别技术（Face++）起家。

印奇曾在 MSRA（微软亚洲研究院）实习，
在那里接触到了深度学习在计算机视觉中的早期应用。
旷视的 Face++ 平台一度是全球最大的人脸识别云服务平台。
公司估值最高时达到约 185 亿元（2025 年胡润全球独角兽榜）。

2026 年 1 月，印奇的"出走"成为行业新闻——
他出任阶跃星辰（StepFun）董事长，
标志着一位 AI 四小龙的创始人正式跨入大模型赛道。
这一人事变动被解读为中国 AI 产业从 CV 时代向 LLM 时代
人才迁移的标志性事件。


% ---- 8.4.3 科大讯飞 ----
\subsection{科大讯飞 (iFlytek) — 语音 AI 的孤勇者}

\subsubsection{刘庆峰 (Liu Qingfeng) — 26 年坚守语音 AI}

\begin{keybox}[人物档案：刘庆峰]
\textbf{生年}：1973 年 \quad \textbf{籍贯}：安徽 \\
\textbf{教育}：中国科学技术大学 \\
\textbf{职务}：科大讯飞创始人、董事长兼 CEO \\
\textbf{成立}：1999 年（中国最早的语音 AI 公司之一）
\end{keybox}

刘庆峰是中国 AI 产业的"长跑运动员"。
1999 年，还在中科大读博的他创立了科大讯飞——
比 AlexNet 早了 13 年，比 GPT 早了 19 年。
在 AI 尚未成为"风口"的年代，他带领讯飞在语音识别领域
艰难前行。科大讯飞的语音合成技术连续多年在国际评测中
名列前茅，讯飞输入法拥有数亿用户。

大模型时代，科大讯飞推出了星火（Spark）大模型。
2025 年，刘庆峰带领讯飞在香港发布了 AI 医疗和智慧课堂产品，
以此为跳板拓展全球市场。他在发布会上表示，
"香港将成为科大讯飞全球化的发射台。"


% ============================================================
%  §8.5  ACADEMIC LEADERS
% ============================================================

\section{学术领袖：象牙塔里的 AI 推动者}
\label{sec:academics}

中国 AI 的学术力量不仅停留在论文层面——
许多顶尖教授直接参与了产业创新，
或者培养了后来创立独角兽公司的学生。
他们是中国 AI 产业的"根系统"。


% ---- 8.5.1 何恺明 ----
\subsection{何恺明 (Kaiming He) — 21 世纪被引最多论文的作者}

\begin{keybox}[人物档案：何恺明]
\textbf{教育}：清华大学物理学学士（2007）、CUHK 计算机博士（2011） \\
\textbf{博士导师}：Xiaoou Tang（汤晓鸥） \\
\textbf{前职}：MSRA（2011--2016）$\to$ FAIR/Meta AI（2016--2024） \\
\textbf{现职}：MIT EECS 终身副教授（2024 年 2 月起）、\\
\quad Google DeepMind 杰出科学家（兼职，2025 年 6 月起） \\
\textbf{代表工作}：ResNet、Mask R-CNN、MAE（Masked Autoencoders）\\
\textbf{荣誉}：AAAI Fellow，未来科学大奖得主
\end{keybox}

何恺明是中国 AI 学者中影响力最大的一位——没有之一。
据 \textit{Nature} 统计，他提出的深度残差网络（ResNet）
是 \textbf{21 世纪被引用次数最多的论文}。
ResNet 提出的残差连接（residual connections）
已成为现代深度学习的标准构件——
从 Transformer（ChatGPT 的核心架构）到 AlphaGo Zero，
从 AlphaFold 到几乎所有的生成式 AI 模型，
残差连接无处不在。

何恺明的人生轨迹本身就是一段传奇：
2003 年广东省高考理科总分第一名（就读于广州执信中学），
进入清华大学物理系；本科毕业后转入 CUHK 计算机系，
师从汤晓鸥攻读博士——他的博士论文是关于单张图像去雾的。
2011--2016 年在 MSRA 工作期间，提出了 ResNet、
Faster R-CNN 等一系列影响深远的工作。
2016 年加入 FAIR（Facebook AI Research），
在 Meta 的八年间产出了 Mask R-CNN、
MAE（Masked Autoencoders）等重要工作。

2024 年 2 月，何恺明加入 MIT EECS 系担任终身副教授。
2025 年 6 月，他更新个人主页，宣布同时兼任 Google DeepMind
的杰出科学家——这一罕见的双重角色反映了学术界和产业界
对他贡献的共同认可。

他目前的研究方向是：构建能够从复杂世界中学习表征和发展智能的
计算机模型。长期目标是用改进的 AI 来增强人类智能。


% ---- 8.5.2 谢赛宁 ----
\subsection{谢赛宁 (Xie Saining) — 从 ResNeXt 到世界模型}

\begin{keybox}[人物档案：谢赛宁]
\textbf{教育}：上海交通大学学士、UCSD 计算机博士 \\
\textbf{前职}：FAIR 研究科学家（4 年）$\to$ Google DeepMind \\
\textbf{现职}：NYU Courant Institute 助理教授、AMI Labs 联合创始人兼首席科学家 \\
\textbf{代表工作}：ResNeXt（与何恺明合著）、ConvNeXt、MoCo（共同作者）
\end{keybox}

谢赛宁是何恺明在 FAIR 期间的核心合作者。
ResNeXt（聚合残差变换）和 ConvNeXt（用现代训练技巧
让纯卷积网络匹敌 Vision Transformer）是他最知名的工作。
MoCo（Momentum Contrast）系列自监督学习论文也有他的贡献。

2022 年，谢赛宁从 FAIR 离职，加入 NYU Courant Institute 担任
助理教授——与图灵奖得主 Yann LeCun 成为同事。
他后来创立了 AMI Labs，作为联合创始人兼首席科学家，
聚焦于世界模型（World Models）——
一种超越纯语言系统、从连续的真实世界传感器经验中学习的 AI。


% ---- 8.5.3 陈丹琦 ----
\subsection{陈丹琦 (Danqi Chen) — NLP 领域的学术新星}

\begin{keybox}[人物档案：陈丹琦]
\textbf{教育}：清华大学学士、Stanford 计算机博士（导师 Christopher Manning） \\
\textbf{现职}：Princeton 大学计算机科学助理教授 \\
\textbf{代表工作}：DrQA（开放域问答先驱）、SimCSE、Dense Passage Retrieval（DPR）\\
\textbf{荣誉}：多项 ACL/EMNLP 最佳论文奖
\end{keybox}

陈丹琦是中国 AI 学者在 NLP 领域的杰出代表。
她在 Stanford 师从 Christopher Manning（NLP 教科书作者）
的博士研究催生了 DrQA——开放域问答系统的开创性工作。
她的 SimCSE（对比学习句子表征）和 DPR（密集段落检索）
是 RAG（检索增强生成）技术路线的重要基础。

在 Princeton，她的实验室同时探索大模型的基础能力和实用技术，
研究方向横跨模型预训练、指令微调、检索增强和推理能力。


% ---- 8.5.4 朱松纯 ----
\subsection{朱松纯 (Zhu Songchun) — 中国 AGI 的学术旗手}

\begin{keybox}[人物档案：朱松纯]
\textbf{生年}：1969 年 \quad \textbf{籍贯}：湖北鄂州 \\
\textbf{教育}：中科大计算机学士、Harvard 计算机博士（导师 David Mumford） \\
\textbf{前职}：UCLA 统计系与计算机系教授（2002--2020）\\
\textbf{现职}：北京通用人工智能研究院（BIGAI）院长、北京大学讲席教授 \& 人工智能研究院/智能学院院长、清华大学基础科学讲席教授 \\
\textbf{荣誉}：IEEE Fellow，全国政协委员，两次 MURI 首席科学家
\end{keybox}

朱松纯是中国 AGI 研究的学术旗手。
他在 Harvard 师从菲尔兹奖得主 David Mumford，
博士研究方向是统计方法在计算机视觉中的应用。
在 UCLA 工作近 20 年间，他建立了视觉、认知和学习实验室（VCLA），
两次担任美国 MURI（多学科大学研究计划）首席科学家——
这是美国国防部资助的最高级别跨学科研究项目。

2020 年，朱松纯做出了一个在学术界引发巨大反响的决定：
\textbf{放弃 UCLA 终身教职，全职回国}。
他在北京筹建了通用人工智能研究院（BIGAI），
同时在北大和清华两所顶尖高校担任讲席教授。

BIGAI 的研究路线与当前主流的"大语言模型"路线有显著差异。
朱松纯主张 AGI 需要\textbf{具有"心"的 AI}——
不仅能处理语言和视觉，还需要理解因果关系、物理直觉和社会认知。
这一立场让他在中国 AI 学术界独树一帜。

2023 年，朱松纯发起了覆盖 9 所高校的
"通用人工智能协同攻关合作体"人才培养计划，
包括北大、清华、浙大、上交大、中科大等——
这是中国在 AGI 人才培养方面最大规模的系统性努力。


% ---- 8.5.5 张钹 ----
\subsection{张钹 (Zhang Bo) — 中国 AI 的学术泰斗}

\begin{keybox}[人物档案：张钹]
\textbf{生年}：1935 年 \quad \textbf{教育}：清华大学 \\
\textbf{职务}：清华大学人工智能研究院名誉院长、中科院院士 \\
\textbf{贡献}：中国 AI 研究的奠基者之一，培养了数代 AI 研究者
\end{keybox}

张钹院士是中国 AI 研究的"活化石"——
在这个领域，他的学术生涯比深度学习的历史还长。
作为清华大学人工智能研究院的奠基人，
他培养了数代中国 AI 研究者，
对中国 AI 学科建设和人才培养的贡献不可估量。
年近九旬仍活跃在学术前沿，多次对 AI 伦理和安全问题发表见解。


% ---- 8.5.6 周志华 ----
\subsection{周志华 (Zhou Zhihua) — 集成学习的世界级权威}

\begin{keybox}[人物档案：周志华]
\textbf{教育}：南京大学计算机科学博士 \\
\textbf{职务}：南京大学人工智能学院院长、计算机系教授 \\
\textbf{代表工作}：深度森林（Deep Forest）、集成学习（Ensemble Learning）\\
\textbf{著作}：《机器学习》（"西瓜书"——中国最畅销的机器学习教科书）\\
\textbf{荣誉}：ACM Fellow、IEEE Fellow、AAAI Fellow、IAPR Fellow、中国 AI 领域引用最高的学者之一
\end{keybox}

周志华是中国机器学习领域的标杆人物。
他在集成学习方面的研究（包括 Boosting、Bagging 的理论分析
和深度森林等创新方法）使他成为该方向的世界级权威。
他撰写的《机器学习》教科书（因封面插图被称为"西瓜书"）
是中国 AI 教育中使用最广泛的教材之一。

2018 年，周志华创立了南京大学人工智能学院——
这是中国 C9 联盟高校中第一个专门的 AI 学院。
他同时是 ACM、IEEE、AAAI、IAPR 四个国际学会的 Fellow——
在中国 AI 学者中极为罕见。


% ---- 8.5.7 马毅 ----
\subsection{马毅 (Ma Yi) — 数学视角的深度学习理论家}

\begin{keybox}[人物档案：马毅]
\textbf{教育}：清华大学自动化学士、UC Berkeley EECS 博士 \\
\textbf{前职}：UIUC 教授 $\to$ Microsoft Research $\to$ UC Berkeley 教授 \\
\textbf{现职}：HKU 计算与数据科学学院院长、UC Berkeley 客座教授 \\
\textbf{代表工作}：稀疏表示人脸识别、鲁棒 PCA、压缩感知 \\
\textbf{荣誉}：IEEE Fellow
\end{keybox}

马毅是一位用严格数学分析深度学习的理论家。
他与 John Wright 等人在稀疏表示（Sparse Representation）
和鲁棒主成分分析（Robust PCA）方面的工作，
是计算机视觉和信号处理领域被引用最多的论文之一。

近年来，马毅的研究聚焦于为深度学习建立数学基础——
试图理解"为什么深度网络能工作"这一根本性问题。
他提出的"压缩-稀疏编码"框架为理解 Transformer 等架构
提供了新的理论视角。

2024 年，马毅出任 HKU（香港大学）新成立的计算与数据科学学院
首任院长，同时保留 UC Berkeley 的客座教授身份——
成为中美 AI 学术合作的重要桥梁。


% ---- 8.5.8 黄铁军 ----
\subsection{黄铁军 (Huang Tiejun) — 类脑计算的探路者}

\begin{keybox}[人物档案：黄铁军]
\textbf{职务}：北京大学计算机学院教授、数字媒体技术研究所所长、\\
\quad 北京智源人工智能研究院（BAAI）院长 \\
\textbf{荣誉}：CCF Fellow、CAAI Fellow、CIE Fellow \\
\textbf{研究方向}：类脑计算、视觉信息处理、数字媒体
\end{keybox}

黄铁军是中国类脑计算（Brain-Inspired Computing）
研究的领军人物之一。他同时担任 BAAI（北京智源人工智能研究院）
院长——BAAI 是中国政府支持的最重要的 AI 研究机构之一。

2025 年 9 月，黄铁军领导的团队发布了 SpikingBrain 1.0——
全球首个基于脉冲神经网络（Spiking Neural Networks）的
大规模 AI 语言模型。与传统 Transformer 架构不同，
脉冲神经网络模仿人脑的工作方式，仅处理相关数据，
能耗远低于传统模型。这一工作代表了中国 AI 在
"非 Transformer"技术路线上的前沿探索。


% ---- 8.5.9 颜水成 ----
\subsection{颜水成 (Yan Shuicheng) — 从学术到产业的多面手}

\begin{keybox}[人物档案：颜水成]
\textbf{教育}：北京大学学士、NUS 博士 \\
\textbf{前职}：NUS 教授 $\to$ 360 集团副总裁 $\to$ Sea AI Lab 负责人 \\
\textbf{现职}：昆仑万维 CEO \\
\textbf{代表工作}：图像超分辨率、视觉注意力机制 \\
\textbf{荣誉}：IEEE Fellow、ACM Fellow、IAPR Fellow
\end{keybox}

颜水成是中国 AI 学者中产学转化最频繁的人物之一。
从 NUS（新加坡国立大学）教授到 360 集团副总裁，
从 Sea AI Lab 负责人到昆仑万维 CEO——
他的职业轨迹跨越了学术界和产业界的多个角色。
这种"教授 $\to$ 大厂高管 $\to$ 上市公司 CEO"的路径，
在中国 AI 领域具有代表性。


% ---- 8.5.10 贾佳亚 ----
\subsection{贾佳亚 (Jia Jiaya) — 从汤晓鸥门生到工业 AI 创业者}

\begin{keybox}[人物档案：贾佳亚]
\textbf{教育}：CUHK 博士 \\
\textbf{前职}：CUHK 教授（长期任职）$\to$ 腾讯优图实验室杰出科学家 \\
\textbf{现职}：HKUST（香港科技大学）讲席教授、冯诺依曼研究所所长；\\
\quad 思谋科技（SmartMore）创始人 \\
\textbf{研究方向}：图像处理、计算摄影、计算机视觉
\end{keybox}

贾佳亚是汤晓鸥在 CUHK 的学术同事和合作者。
在 CUHK 任教期间，他在图像去模糊、超分辨率等方向
产出了一系列高影响力的工作。

2019 年，贾佳亚创立了思谋科技（SmartMore），
专注于工业视觉 AI——将计算机视觉技术应用于制造业
的质量检测和智能化升级。到 2025 年，思谋科技的年收入
突破 1.5 亿美元。2026 年 3 月，思谋正式向港交所提交
IPO 申请，目标成为"工业 AI Agent 第一股"。

贾佳亚的故事展示了中国 AI 学者的另一条成功路径：
不是做通用大模型，而是\textbf{将 AI 技术深度嵌入垂直行业}。


% ============================================================
%  §8.6  AI CHIP PIONEERS
% ============================================================

\section{AI 芯片先驱：硅基战场的中国力量}
\label{sec:chips}

在美国对华芯片出口管制的大背景下，
中国 AI 芯片公司成为了地缘政治博弈中的关键角色。
它们面临的挑战是双重的：既要追赶 NVIDIA 的技术领先，
又要在受限的制造工艺下实现商业化。


% ---- 8.6.1 寒武纪 ----
\subsection{寒武纪 (Cambricon) — 天才兄弟的"中国 NVIDIA"梦}

\subsubsection{陈天石 (Chen Tianshi) \& 陈云霁 (Chen Yunji) — 天才家族}

\begin{keybox}[人物档案：陈天石 \& 陈云霁]
\textbf{陈天石}：1985 年生，江西南昌。中科大少年班（16 岁入学），中科院计算所博士。寒武纪创始人、董事长兼 CEO。\\
\textbf{陈云霁}：陈天石之兄。14 岁考入中科大少年班，23 岁获中科院计算所博士。中科院计算所研究员，寒武纪联合创始人。
\end{keybox}

陈氏兄弟的故事堪称中国 AI 芯片领域最具传奇色彩的创业叙事。
哥哥陈云霁 14 岁考入中科大少年班，
弟弟陈天石 16 岁进入同一所学校的同一个班——
这个家庭被外界称为"天才之家"。

兄弟二人毕业后在中科院计算所被破格录用，
成为龙芯团队最年轻的成员。2010 年前后，
他们开始在团队内部提出"AI 芯片"的研发构想——
主张通过硬件架构创新来加速深度学习计算。

2016 年，31 岁的陈天石创立寒武纪，选择了\textbf{全栈路线}——
从芯片架构到编程框架到软件工具链的完整技术栈。
这条路线意味着更长的研发周期和更大的资金消耗，
但也意味着更强的生态控制力。

寒武纪的历程充满了戏剧性反转：
\begin{itemize}[noitemsep]
  \item \textbf{2020 年}：科创板上市，但招股书披露三年累计亏损 16 亿元，
        市场一片质疑。
  \item \textbf{2024 年}：股价开始暴涨，全年涨幅近 10 倍。
  \item \textbf{2025 年}：首次实现年度盈利——营收超 60 亿元（同比增长 4 倍以上），
        净利润超 18.5 亿元。股价一度突破 1595 元/股，
        \textbf{超越贵州茅台成为 A 股新一代"股王"}。
  \item 陈天石的持股财富一度达到 \textbf{1902 亿元}。
\end{itemize}

\begin{corebox}[寒武纪的"十年穿越无人区"]
寒武纪的成功验证了一个被许多投资者低估的事实：
\textbf{AI 芯片是一场十年级别的长跑}。

陈天石从 2016 年创业到 2025 年首次盈利，用了整整 9 年。
在这 9 年中，公司累计亏损数十亿元，多次面临市场质疑和资金压力。
但他坚持了全栈路线——而不是选择更容易短期变现的 IP 授权模式。

2025 年的盈利转折点来自两个因素：
(1) DeepSeek 效应带动了全中国 AI 算力需求的爆发式增长；
(2) 美国芯片出口管制使华为昇腾和寒武纪成为国内唯二的可选方案。
陈天石"10 年坐冷板凳"的故事，
与梁文锋"从量化基金到 AGI"的跨界叙事一起，
构成了中国 AI 芯片产业最动人的创业传奇。
\end{corebox}


% ---- 8.6.2 摩尔线程 ----
\subsection{摩尔线程 (Moore Threads) — NVIDIA 中国前掌门的逆袭}

\subsubsection{张建中 (Zhang Jianzhong / James Zhang)}

\begin{keybox}[人物档案：张建中]
\textbf{生年}：$\sim$1965 年 \\
\textbf{教育}：南京理工大学计算机科学 \\
\textbf{前职}：HP $\to$ Dell $\to$ NVIDIA 中国区副总裁兼总经理（14 年）\\
\textbf{职务}：摩尔线程创始人、董事长兼 CEO \\
\textbf{财富}：$\sim$43 亿美元（上市后）
\end{keybox}

张建中在 NVIDIA 工作了 14 年，
从 2005 年加入到 2019 年离职，
升任副总裁兼中国区总经理——
他是 NVIDIA 在中国业务扩张最关键时期的操盘手。

2020 年，张建中创立摩尔线程，
目标是打造"中国的 NVIDIA"——
一家同时覆盖 GPU 图形渲染和 AI 加速的全功能芯片公司。
公司名取自"摩尔定律"（Moore's Law）——
暗示了对芯片性能持续提升的信念。

2025 年 12 月 5 日，摩尔线程在上海科创板上市，
首日暴涨超过 400\%，创造了中国 AI 芯片 IPO 的纪录。
张建中的持股财富因此达到约 43 亿美元——
从 NVIDIA 高管到十亿美元级别的创始人，
这一"从打工者到老板"的故事在中国科技圈广为流传。

\begin{whybox}[从 NVIDIA 出走的中国 GPU 军团]
张建中并非唯一离开 NVIDIA 在中国创业的高管。
"GPU 四小龙"——摩尔线程、沐曦（Metax）、
燧原科技（Enflame）、壁仞科技（Biren）——
中有多位核心创始人具有 NVIDIA 或其他国际芯片公司的背景。

这一现象反映了 AI 芯片产业人才流动的独特规律：
与软件不同，芯片设计需要 10--20 年的经验积累。
中国 AI 芯片创业公司的技术核心几乎全部来自
在国际芯片巨头（NVIDIA、AMD、Qualcomm）工作多年的资深工程师。
美国的芯片出口管制在限制中国获取先进芯片的同时，
也为这些回流人才提供了前所未有的创业机遇。
\end{whybox}


% ---- 8.6.3 壁仞 ----
\subsection{壁仞科技 (Biren Technology) — 港交所 2026 开年第一股}

壁仞科技成立于 2019 年，总部位于上海。
2026 年 1 月 2 日，壁仞在港交所挂牌上市——
成为 2026 年中国资本市场的第一个 IPO。
与摩尔线程和寒武纪一起，壁仞的上市标志着
中国 AI 芯片公司正在集体进入公开资本市场。

壁仞的创始团队汇聚了来自 AMD、NVIDIA、S3 Graphics
等国际芯片公司的资深专家。
公司最初专注于通用 GPU 芯片，
后在 AI 训练和推理加速方向加大投入——
其 BR100 系列芯片对标 NVIDIA A100。


% ============================================================
%  §8.7  EMBODIED AI / ROBOTICS
% ============================================================

\section{具身智能：中国机器人创业浪潮}
\label{sec:china-robotics}

2025--2026 年，中国具身智能（Embodied AI）赛道迎来了前所未有的爆发。
从春晚舞台上翻跟头的人形机器人，到工厂流水线上搬运物料的双足机器人，
从 B 站百万粉丝的"稚晖君"到北大教授下海创业——
中国正在以惊人的速度成为全球人形机器人产业的中心之一。

这一浪潮的背后有三股力量：
(1) 大模型技术为机器人的感知和决策提供了全新范式（VLA 模型）；
(2) 中国强大的制造业供应链极大降低了机器人硬件成本；
(3) 人口老龄化和劳动力短缺创造了刚性需求。

\begin{corebox}[中国机器人创业公司一览（2026 年 3 月）]
\begin{tabular}{lllr}
\textbf{公司} & \textbf{创始人} & \textbf{方向} & \textbf{估值} \\
\hline
宇树科技 Unitree & 王兴兴 & 四足/人形机器人 & $\sim$80 亿元 \\
智元机器人 Agibot & 彭志辉 & 人形机器人 & $\sim$21 亿美元 \\
银河通用 Galbot & 王鹤 & 通用机器人 & $\sim$30 亿美元 \\
傅利叶智能 Fourier & 顾捷 & 康复/人形机器人 & $\sim$10 亿美元 \\
\end{tabular}
\end{corebox}


% ---- 8.7.1 宇树科技 ----
\subsection{宇树科技 (Unitree Robotics) — 春晚上的机器人帝国}

\subsubsection{王兴兴 (Wang Xingxing) — 从 200 块钱造机器人到全球六成市场份额}

\begin{keybox}[人物档案：王兴兴]
\textbf{生年}：1990 年 \quad \textbf{籍贯}：浙江宁波余姚 \\
\textbf{教育}：浙江理工大学机电工程学士、上海大学机电工程硕士 \\
\textbf{职务}：宇树科技创始人、CEO 兼 CTO \\
\textbf{成就}：四足机器人全球市场份额超 60\%，2025 年春晚机器人表演
\end{keybox}

王兴兴的创业故事始于一个 200 元人民币的小型双足机器人。

1990 年出生于浙江宁波余姚一个普通家庭，
王兴兴从小就对机械和电子充满痴迷。小学时他的第一个发明是一辆
风力小车，老师和同学们称他为"爱折腾的小孩"。
他在学业上并不以传统标准衡量出色——英语成绩一度不理想——
但动手能力极强。

2009 年进入浙江理工大学机电工程专业后，大一就用约 200 元
零件手工打造了他的第一个双足机器人。
2013 年进入上海大学攻读硕士，2015 年用约两万元开发了一款
名为 XDog 的四足机器人原型，作为毕业论文项目。
这个项目成为宇树科技的技术原点。

2016 年硕士毕业后，王兴兴短暂加入大疆创新（DJI），但仅两个月后便辞职，
创立了宇树科技。"宇树"这个名字的含义是
"在广阔宇宙中做一棵枝繁叶茂的科技树"。
公司头三年极为艰难——据王兴兴回忆，
有时候连发工资都困难。2018 年底，宇树开始产生收入并获得新一轮融资。

宇树的爆发始于两个关键节点：
\begin{itemize}[noitemsep]
  \item \textbf{2024 年}：宇树 Go2 机器狗以 1600 美元的起售价震惊行业——
        这一价格仅为 Boston Dynamics Spot（74500 美元）的 2\%。
        四足机器人全球市场份额突破 60\%。
  \item \textbf{2025 年春晚}：宇树的人形机器人和四足机器狗在
        中央电视台春节联欢晚会上表演武术和秧歌，
        超过 7 亿观众观看。这一事件让宇树从科技圈"出圈"，
        成为中国家喻户晓的机器人品牌。
\end{itemize}

2026 年 3 月，王兴兴在公开演讲中预测：
"今年年中，全球尤其是中国的机器人有望在百米冲刺中跑进 10 秒以内，
超越博尔特 9.58 秒的世界纪录。"

\begin{whybox}[宇树为什么能把机器人做到"白菜价"？]
王兴兴在多次采访中解释了宇树的成本优势来源：
\begin{enumerate}[noitemsep]
  \item \textbf{自研核心零部件}：电机、减速器、控制器全部自主研发，
        避免了对昂贵进口组件的依赖。
  \item \textbf{中国制造业供应链}：深圳和长三角的电子元器件和
        精密制造生态系统提供了全球最低成本的零部件采购渠道。
  \item \textbf{消费级硬件思维}：王兴兴的理念是"用消费电子的方法做机器人"——
        追求极致的成本优化和大规模量产，而非传统机器人行业的小批量高定价。
\end{enumerate}
这一策略与 DeepSeek 用软件效率弥补硬件限制的思路异曲同工——
\textbf{在资源约束下寻找非对称优势}。
\end{whybox}


% ---- 8.7.2 智元机器人 ----
\subsection{智元机器人 (Agibot) — B 站 UP 主的钢铁侠梦想}

\subsubsection{彭志辉 (Peng Zhihui / 稚晖君) — 从华为天才少年到机器人创业者}

\begin{keybox}[人物档案：彭志辉]
\textbf{生年}：1993 年 \quad \textbf{籍贯}：江西吉安 \\
\textbf{教育}：电子科技大学生命科学学士、信息与通信工程硕士 \\
\textbf{前职}：OPPO AI 实验室算法工程师 $\to$ 华为"天才少年"计划 \\
\textbf{职务}：智元机器人联合创始人、总裁兼 CTO \\
\textbf{别名}：稚晖君、"野生钢铁侠" \\
\textbf{社交影响力}：B 站 284 万粉丝、微博 103 万粉丝
\end{keybox}

彭志辉是中国 AI 创业者中最独特的一位——
他的"出圈"不是靠融资或产品发布，而是靠 B 站的硬核 DIY 视频。

1993 年出生于江西吉安永新县一个农村家庭，
彭志辉从小就对电子设备着迷。别的孩子在打游戏，
他却在研究游戏机的运行原理，最终直接把游戏机拆了。
2015 年本科毕业于电子科技大学生命科学与技术学院——
这个专业与他后来的 AI 和机器人事业看似毫无关联，
但恰恰体现了他"跨界思维"的特质。
2018 年硕士毕业于电子科技大学信息与通信工程学院，
在校期间是竞赛狂人：全国大学生电子设计竞赛一等奖、
飞思卡尔智能车竞赛赛区一等奖、挑战杯金奖、
微软极限编程竞赛一等奖。

2017 年开始，彭志辉以"稚晖君"为网名在 B 站发布硬核 DIY 视频。
他制作的项目令人叹为观止：
\begin{itemize}[noitemsep]
  \item 一台不比一元硬币大的迷你 Linux 电脑
  \item 能缝合葡萄皮的机械臂"Dummy"
  \item 自动驾驶自行车
  \item 自平衡机器人
\end{itemize}
这些视频让他在 B 站迅速走红，被粉丝称为"野生钢铁侠"。
2022 年获得 B 站年度百大 UP 主称号。

硕士毕业后，彭志辉以 SSP（Super Special Offer）入职 OPPO 研究院 AI 实验室。
2020 年，他通过华为"天才少年"计划加入华为，
从事昇腾 AI 芯片和 AI 算法研究。
华为天才少年计划的年薪高达 100--200 万元，
是华为招揽顶尖年轻人才的旗舰项目。

2022 年 12 月 27 日，彭志辉在微博发文宣布离开华为，开始创业。
2023 年 2 月，他联合创立了智元机器人（Agibot），担任 CTO。

智元机器人的发展速度令人惊叹：
\begin{itemize}[noitemsep]
  \item \textbf{2023 年 8 月}：首款具身智能机器人"远征 A1"公开亮相。
  \item \textbf{2025 年 3 月}：发布 GO-1（Genie Operator-1），
        号称"全球首个通用具身基座大模型"——
        一个能让机器人理解自然语言指令并执行物理世界任务的 VLA 模型。
  \item \textbf{估值}：截至 2025 年 3 月达到 21 亿美元，投资方包括 BYD 和高瓴。
\end{itemize}

\begin{whybox}[从 B 站 UP 主到机器人公司 CTO：彭志辉的"反传统"路径]
彭志辉的职业路径颠覆了中国 AI 创业的常规模式。
大多数六小虎创始人的路径是"清华/北大 $\to$ 海外博士 $\to$ 大厂 $\to$ 创业"；
而彭志辉的路径是"非顶尖高校 $\to$ B 站走红 $\to$ 华为天才少年 $\to$ 创业"。

他的成功揭示了一个被低估的趋势：
\textbf{在社交媒体时代，技术影响力不再完全取决于学术出身}。
B 站的 284 万粉丝不仅为他提供了品牌效应，
更重要的是建立了一个庞大的技术人才关注网络——
这在招聘和融资时都转化为了巨大的优势。
\end{whybox}


% ---- 8.7.3 银河通用 ----
\subsection{银河通用 (Galbot) — 北大教授的通用机器人之路}

\subsubsection{王鹤 (Wang He) — 从仿真到现实的具身智能先锋}

\begin{keybox}[人物档案：王鹤]
\textbf{职务}：银河通用机器人创始人兼 CTO、北京大学助理教授 \\
\textbf{研究方向}：具身智能、机器人仿真、感知与决策 \\
\textbf{估值}：2025 年底完成 3 亿美元融资，估值约 30 亿美元
\end{keybox}

王鹤的创业切入点独树一帜：\textbf{仿真}（Simulation）。
在大多数机器人公司专注于硬件本体或算法时，
王鹤认为仿真是打通具身智能"最后一公里"的关键——
在虚拟环境中让机器人学会感知和决策，
再将能力迁移到真实物理世界。

银河通用成立于 2023 年，发展速度极快：
2025 年底完成 3 亿美元融资，估值达到 30 亿美元，
投资方包括宁德时代（CATL）、中金资本、GGV Capital、
中科院基金等——这一阵容反映了从新能源到国家资本
对具身智能赛道的集体看好。

2026 年央视春晚上，银河通用的机器人也参与了表演，
进一步提升了品牌知名度。


% ---- 8.7.4 傅利叶智能 ----
\subsection{傅利叶智能 (Fourier Intelligence) — 从康复机器人到通用人形}

\subsubsection{顾捷 (Gu Jie) — 十年康复积累，一朝转型人形}

\begin{keybox}[人物档案：顾捷]
\textbf{职务}：傅利叶智能创始人兼 CEO \\
\textbf{背景}：上海交通大学，康复医疗机器人领域深耕十余年 \\
\textbf{产品}：GR-1（全球首款量产通用人形机器人）、GR-2、GR-3
\end{keybox}

傅利叶智能的差异化在于它的\textbf{医疗基因}。
公司成立于 2015 年，最初专注于康复医疗机器人——
其产品已部署在全球 2000+ 家医疗机构。
这十年的康复机器人经验为傅利叶提供了三个核心优势：
精密电机控制技术、人体运动学理解、以及医疗级安全标准。

2023 年，傅利叶发布 GR-1——被称为全球首款量产通用人形机器人。
2024 年发布 GR-2（175cm、63kg、53 个自由度），
2026 年发布 GR-3 Care-Bot，专注于护理和情感交互。

顾捷在接受 36Kr 采访时指出：
"大模型是让人形机器人从实验室走向量产的关键催化剂。
但真正的壁垒不在算法，而在硬件的可靠性和供应链管理——
这恰恰是中国汽车和消费电子产业积累了几十年的能力。"

\begin{corebox}[中国机器人的"iPhone 时刻"？]
2025--2026 年的中国机器人创业浪潮，
与 2007--2012 年的智能手机革命有惊人的相似之处：
\begin{itemize}[noitemsep]
  \item \textbf{大模型 = iOS/Android}：VLA（Vision-Language-Action）模型
        为机器人提供了统一的"操作系统"，
        就像 iOS 和 Android 为手机生态提供了标准化平台。
  \item \textbf{中国供应链 = 深圳制造}：如同智能手机的硬件成本
        因深圳供应链而急剧下降，机器人的零部件成本
        也在中国制造业的规模效应下快速降低。
  \item \textbf{消费级定价}：宇树 Go2（1600 美元）之于机器狗，
        正如小米手机之于智能手机——用极致性价比打开大众市场。
\end{itemize}
但也有关键差异：机器人需要在物理世界中安全运作，
法律、伦理和安全标准的建立远比软件生态复杂。
王兴兴自己也坦言："家用机器人还面临法律、道德、伦理等问题，
工业和商业应用会更快落地。"
\end{corebox}


% ============================================================
%  §8.8  AUTONOMOUS DRIVING
% ============================================================

\section{自动驾驶：智能驾驶的中国力量}
\label{sec:autonomous-driving}

如果说大模型竞赛是 2023--2025 年中国 AI 最热的赛道，
那么自动驾驶则是投入最深、周期最长的 AI 应用方向。
从 2016 年至今，中国自动驾驶公司经历了技术验证、商业化探索、
上市潮和整合期四个阶段。2024--2025 年，
多家中国自动驾驶公司密集登陆美股和港股，
标志着这个赛道从"讲故事"进入"算账"阶段。


% ---- 8.8.1 小马智行 ----
\subsection{小马智行 (Pony.ai) — 中国自动驾驶的"楼教主"传说}

\subsubsection{彭军 (James Peng) — 从百度到 Pony.ai}

\begin{keybox}[人物档案：彭军]
\textbf{前职}：百度自动驾驶事业部首席架构师、Google 工程师 \\
\textbf{职务}：小马智行联合创始人兼 CEO \\
\textbf{上市}：2024 年 11 月 NASDAQ（PONY）、2025 年 11 月港交所双重上市 \\
\textbf{港股 IPO}：募资最高 77 亿港元，2025 年全球自动驾驶最大 IPO
\end{keybox}

\subsubsection{楼天城 (Lou Tiancheng / ACRush) — 传说中的"楼教主"}

\begin{keybox}[人物档案：楼天城]
\textbf{别名}：楼教主、ACRush \\
\textbf{教育}：清华大学计算机系学士、清华大学姚期智班博士 \\
\textbf{博士导师}：姚期智（图灵奖得主）\\
\textbf{前职}：Google $\to$ Google X 自动驾驶 $\to$ Quora $\to$ 百度自动驾驶 \\
\textbf{职务}：小马智行联合创始人兼 CTO \\
\textbf{竞赛成绩}：Google Code Jam 两届冠军、TopCoder 十年奖牌获得者、\\
\quad IOI 铜牌（2004）、ICPC 世界总决赛亚军两次
\end{keybox}

在中国编程竞赛圈，"楼教主"三个字是一个传说。

楼天城是中国竞赛编程历史上最具统治力的选手之一。
他的 TopCoder ID "ACRush"在全球编程竞赛社区家喻户晓——
连续十年在 TopCoder Open 中获得奖牌，
两次夺得 Google Code Jam 全球冠军（2008、2009），
代表清华大学两次获得 ICPC 世界总决赛亚军。
2008 年被 MIT Technology Review 评为 35 岁以下创新者。

他在清华大学姚期智班（Institute for Theoretical Computer Science）
完成了从本科到博士的全部学业——导师正是图灵奖得主姚期智。
博士毕业后加入 Google，2015 年转入 Google X 自动驾驶团队，
后短暂加入 Quora，再被 Andrew Ng 招入百度自动驾驶部门——
成为百度最年轻的 T10 级工程师。

2016 年底，楼天城与彭军（时任百度自动驾驶首席架构师）共同离开百度，
在硅谷创立了小马智行。公司名中的"Pony"既是"小马"的英文，
也暗合了马化腾（Pony Ma）——小马智行的早期投资者之一。

小马智行的上市之路堪称中国自动驾驶行业的里程碑：
\begin{itemize}[noitemsep]
  \item \textbf{2024 年 11 月}：以"PONY"为代码在 NASDAQ 挂牌上市，
        募资约 4.52 亿美元，创当年美股自动驾驶最大 IPO。
  \item \textbf{2025 年 11 月}：以"2026"为代码在港交所完成双重上市，
        募资最高 77 亿港元——成为 2025 年全球自动驾驶最大 IPO，
        也是 2022 年以来全球自动驾驶行业规模最大的 IPO。
  \item ARK Invest（方舟投资）、Wellington、Fidelity 等国际顶级机构积极买入。
\end{itemize}


% ---- 8.8.2 文远知行 ----
\subsection{文远知行 (WeRide) — 全球首家双重上市的 Robotaxi 公司}

\begin{keybox}[人物档案：韩旭 (Han Xu)]
\textbf{教育}：University of Illinois at Urbana-Champaign 博士 \\
\textbf{前职}：Missouri S\&T 教授 \\
\textbf{职务}：文远知行创始人兼 CEO \\
\textbf{上市}：2024 年 10 月 NASDAQ（WRD）、2025 年 11 月港交所（0800）
\end{keybox}

文远知行是全球第一家同时在美国和香港两地上市的 Robotaxi 公司。
2024 年 10 月在 NASDAQ 挂牌，成为全球首家上市的通用自动驾驶技术公司。
2025 年 11 月在港交所完成双重主要上市，
募资约 23.9 亿港元——创下港交所 18C 章节
（专为特专科技公司设计的上市机制）最大规模 IPO 纪录。

文远知行的业务覆盖 Robotaxi、Robobus、Robotruck 和自主泊车等多个场景，
是中国自动驾驶赛道中"全场景"布局最广的公司之一。


% ---- 8.8.3 Momenta ----
\subsection{Momenta — 量产级自动驾驶的"第三极"}

\begin{keybox}[人物档案：曹旭东 (Cao Xudong)]
\textbf{教育}：清华大学计算机系 \\
\textbf{前职}：Microsoft Research Asia（MSRA）\\
\textbf{职务}：Momenta 创始人兼 CEO \\
\textbf{估值}：2025 年融资估值约 60 亿美元
\end{keybox}

在中国自动驾驶市场，形成了"自研、华为、Momenta"三足鼎立的格局——
"自研"指车企自主开发的系统，"华为"指华为智能驾驶方案，
而 Momenta 则代表了第三方技术供应商的最高水平。

曹旭东毕业于清华大学计算机系，曾任职于 MSRA。
2016 年创立 Momenta，走了一条与小马智行截然不同的路线：
\textbf{不做 L4 Robotaxi，而是做量产级 L2+/L3 辅助驾驶}。
这一策略的核心逻辑是"用量产数据飞轮加速技术迭代"——
通过与车企合作，在量产车上部署系统，收集海量真实驾驶数据，
再用数据反哺算法提升。

Momenta 的客户名单极为亮眼：BMW、Mercedes-Benz、SAIC、丰田等
全球顶级车企都采用了 Momenta 的智能驾驶方案。
投资方包括通用汽车、腾讯、Temasek 和丰田。
2025 年估值达到约 60 亿美元，并于 2026 年 3 月向港交所提交
IPO 申请，计划募资 10 亿美元。


% ---- 8.8.4 蔚小理 ----
\subsection{蔚小理 (NIO、XPeng、Li Auto) — 新势力的 AI 雄心}

中国新能源汽车"三驾马车"——蔚来、小鹏、理想——
在智能驾驶方向的投入和野心远超传统车企。

\paragraph{何小鹏 (He Xiaopeng) — 从 UC 浏览器到"Physical AI"}
小鹏汽车（XPeng）创始人何小鹏是中国科技创业者中少有的
"二次创业"成功案例。他创立的 UC 浏览器被阿里巴巴以约 40 亿美元收购后，
2014 年创立小鹏汽车。何小鹏将 AI 定义为小鹏的核心差异化——
公司已训练了百亿参数级别的云端驾驶模型，
使用超过一亿公里的真实驾驶数据进行预训练。
2026 年 2 月，何小鹏在全员信中提出"Physical AI"战略，
目标是成为全球首家在同一年内实现人形机器人、飞行汽车和
Robotaxi 三大前沿 AI 业务量产的科技公司。

\paragraph{李斌 (Li Bin) — 蔚来的自动驾驶芯片野心}
蔚来创始人李斌走的是"全栈自研"路线——
2025 年发布自研智能驾驶芯片"神玑 NX9031"（5nm 工艺），
成为中国首家同时自研车机系统芯片和智能驾驶芯片的车企。

\paragraph{李想 (Li Xiang) — 理想的"端到端"路线}
理想汽车创始人李想是汽车科技媒体"汽车之家"的创始人（后被收购），
二次创业的理想汽车以"增程式"技术路线在中国新能源市场独树一帜。
在智能驾驶方面，理想采用"端到端"大模型方案，
并在 2025 年实现了 L3 级别辅助驾驶的大规模推送。

\begin{whybox}[自动驾驶为什么是中国 AI 的"长征"？]
中国自动驾驶赛道从 2016 年至今已走过近十年。
与大模型赛道的"短期爆发"不同，自动驾驶是一场需要
持续投入、长期验证、逐步商业化的漫长征途。
小马智行从创立到上市用了 8 年，文远知行同样如此。
这条赛道的生存法则是：
\textbf{活得够久，比跑得够快更重要}。
2024--2025 年的密集上市潮，标志着第一批"幸存者"
开始进入公开资本市场——但距离真正的大规模盈利，
可能还需要另一个五年。
\end{whybox}


% ============================================================
%  §8.9  OPEN SOURCE ECOSYSTEM & AI INFRA
% ============================================================

\section{中国开源生态与 AI 基础设施}
\label{sec:open-source}

DeepSeek-V3/R1 的开源发布不仅震动了全球 AI 产业，
更深刻地改变了中国 AI 开源生态的格局。
据 Stanford HAI 2025 年分析报告，至少 12 家中国机构
在生产 SOTA 级别的开源/开放权重 AI 模型，
中国在全球开放权重 LLM 的分发量上已超越美国。

但中国的开源生态远不止 DeepSeek 和 Qwen——
从分布式训练框架到推理加速平台，从端侧小模型到学术开源工具，
一个多层次的开源技术栈正在快速成型。


% ---- 8.9.1 开源模型家族 ----
\subsection{开源模型四大家族}

\paragraph{DeepSeek 开源系列}
DeepSeek-V3/R1 的开放权重发布是 2025 年全球 AI 最具影响力的事件之一。
R1 论文发表于 \textit{Nature}，模型权重以 MIT 许可证开源——
这意味着任何人、任何组织都可以自由使用、修改和商用。
这一决策的影响是链式反应式的：
全球数千家公司和研究机构基于 DeepSeek 模型构建了自己的产品，
零一万物等公司直接转向基于 DeepSeek 开源模型的应用层路线。

\paragraph{Qwen 开源系列（阿里巴巴）}
阿里通义千问（Qwen）团队是中国开源大模型生态中最活跃的贡献者。
截至 2026 年初，Qwen 系列累计开源模型超过 400 个，
在 Hugging Face 上的衍生模型突破 20 万个，总下载量超过 10 亿次——
\textbf{超越 Meta 的 Llama 系列成为全球下载量最高的开源模型家族}。
Qwen 的成功在于其"全系列"策略：
从 0.5B 到 110B 参数，从纯文本到多模态，
从基座模型到对话模型，覆盖了几乎所有使用场景。

\paragraph{ChatGLM / GLM 开源系列（智谱/清华 KEG）}
智谱 AI 背后的清华 KEG 实验室是中国最早推动大模型开源的学术力量。
ChatGLM 系列从 2023 年发布以来，一直是中文开源社区最受欢迎的模型之一。
GLM 的技术路线（自回归填空预训练目标）与 GPT 和 BERT 都不同，
为中文 NLP 社区提供了独特的技术选择。

\paragraph{InternLM 开源系列（上海 AI 实验室）}
上海人工智能实验室（Shanghai AI Lab）发布的 InternLM 系列
是中国国家级 AI 研究机构在开源领域的代表作。
实验室由汤晓鸥创立，其去世后由\textbf{周伯文 (Zhou Bowen)}
接任院长——周伯文曾任 IBM Watson AI 首席科学家和
京东集团副总裁兼 AI 研究院院长。
InternLM 系列的特色是"强推理+长上下文"，
配合 InternVL（视觉-语言模型）构成了完整的多模态开源生态。


% ---- 8.9.2 ColossalAI ----
\subsection{ColossalAI / 潞晨科技 — 让大模型训练"人人可及"}

\subsubsection{尤洋 (Yang You) — 从 UC Berkeley 到分布式训练先锋}

\begin{keybox}[人物档案：尤洋]
\textbf{教育}：清华大学学士、UC Berkeley 计算机科学博士 \\
\textbf{现职}：新加坡国立大学 (NUS) Presidential Young Professor、\\
\quad 潞晨科技（HPC-AI Tech）创始人兼 CEO \\
\textbf{代表工作}：LARS/LAMB 大批量优化器、Colossal-AI 分布式训练框架 \\
\textbf{荣誉}：ACM Gordon Bell Prize 提名
\end{keybox}

尤洋是全球分布式 AI 训练领域最重要的华人学者之一。
他在 UC Berkeley 的博士研究聚焦于大规模深度学习训练优化——
LARS（Layer-wise Adaptive Rate Scaling）和 LAMB（Layer-wise Adaptive
Moments optimizer for Batch training）是他的标志性工作，
被 Google、NVIDIA 等公司广泛采用来加速大模型训练。

2021 年，尤洋创立潞晨科技（HPC-AI Tech），
开发了 Colossal-AI 开源分布式训练框架。
这个项目在 GitHub 上获得超过 41000 Star，190+ 贡献者，
已发布 50 个版本——是全球最活跃的 AI 训练基础设施开源项目之一。

Colossal-AI 的核心价值在于\textbf{让大模型训练民主化}：
通过自动并行化、内存优化和通信加速，
使中小团队也能在有限 GPU 资源上训练大规模模型。
其用户包括 AWS、Meta、BioMap 等国际机构。


% ---- 8.9.3 硅基流动 ----
\subsection{硅基流动 (SiliconFlow) — AI 推理的"水电煤"}

\subsubsection{袁进辉 (Yuan Jinhui) — 从 MSRA 到推理加速创业}

\begin{keybox}[人物档案：袁进辉]
\textbf{教育}：清华大学计算机系博士 \\
\textbf{前职}：Microsoft Research Asia \\
\textbf{职务}：硅基流动创始人兼 CEO \\
\textbf{融资}：Pre-A 轮亿元人民币（华创资本领投），美团战略投资
\end{keybox}

如果说 DeepSeek 证明了中国团队能以低成本\textbf{训练}世界级模型，
那么硅基流动要解决的是下一个问题：
如何以低成本\textbf{部署和推理}这些模型。

硅基流动成立于 2023 年 8 月，创始人袁进辉是清华大学计算机系博士，
曾在 MSRA 任职多年。公司的核心产品 SiliconCloud 是一个
一站式 AI 云服务平台，提供从模型训练、微调到推理部署的全链路服务。

SiliconCloud 上线不到一年即实现爆发式增长：
平台用户数突破 300 万，日均调用量超过千亿 Token。
其自研推理引擎 SiliconLLM 通过软硬件联合优化，
大幅降低推理成本——并率先实现了在华为昇腾芯片上
部署满血版 DeepSeek-R1/V3，获得持平 NVIDIA 高端 GPU 的效果。

硅基流动的战略定位是"AI 时代的水电煤"——
不做模型本身，而是为所有模型开发者和应用方
提供极致高效、低成本的推理基础设施。


% ---- 8.9.4 面壁智能 ----
\subsection{面壁智能 (ModelBest) — 端侧 AI 的学术创业样本}

\begin{keybox}[面壁智能核心信息]
\textbf{创始团队}：刘知远（清华大学 NLP 教授）等清华 NLP 学者 \\
\textbf{代表产品}：MiniCPM 系列——高性能端侧小模型 \\
\textbf{最新发布}：MiniCPM-o4.5（全球首个全双工全模态 AI 模型，2026 年 2 月）\\
\textbf{路线}：端侧部署优先、开源优先
\end{keybox}

面壁智能是清华大学 NLP 实验室的学术创业代表。
其核心技术来源于清华 NLP 组的 CPM（Chinese Pre-trained Models）系列——
中国最早的大规模中文预训练模型之一。

面壁智能的差异化在于\textbf{端侧小模型}：
MiniCPM 系列以仅 9B 参数实现了接近 GPT-4o 级别的多模态能力，
可以直接在手机等端侧设备上运行。
2026 年 2 月发布的 MiniCPM-o4.5 更进一步——
这是全球首个全双工（Full-Duplex）全模态 AI 模型，
能够"边看、边听、边说"，无需等待用户说完即可实时响应。

这一路线的战略意义在于：
\textbf{当算力无法集中时，让 AI 分布到每一台设备上}。


% ---- 8.9.5 更多学术力量 ----
\subsection{更多学术开源力量}

\subsubsection{邱锡鹏 (Qiu Xipeng) — MOSS 与复旦 NLP}

\begin{keybox}[人物档案：邱锡鹏]
\textbf{教育}：复旦大学理学学士、博士 \\
\textbf{职务}：复旦大学教授、上海创智学院 OpenMOSS 团队负责人 \\
\textbf{代表工作}：MOSS 大语言模型、FudanNLP 工具包、FastNLP \\
\textbf{著作}：《神经网络与深度学习》（蒲公英书）\\
\textbf{荣誉}：CAAI Fellow、ACL 2017/2024 杰出论文奖、CCF-ACM AI 青年奖
\end{keybox}

邱锡鹏是中国 NLP 开源社区最重要的学术推动者之一。
2023 年 2 月，他带领复旦团队发布了 MOSS——
中国第一个公开发布的对话大语言模型（16B 参数），
比百度文心一言和阿里通义千问的发布都更早。
虽然 MOSS 在性能上很快被后来者超越，
但其"开源优先"的理念和快速公开的做法，
为中国大模型开源生态树立了重要先例。

邱锡鹏撰写的教科书《神经网络与深度学习》（被称为"蒲公英书"）
是中文深度学习教材中使用最广泛的之一，
与周志华的"西瓜书"一起构成了中国 AI 教育的两大支柱。

\begin{corebox}[中国开源 AI 生态的全球影响]
据 MIT Technology Review 2026 年 2 月报道和 Stanford HAI 分析：
\begin{itemize}[noitemsep]
  \item 中国开放权重模型在 Hugging Face 上的累计下载量
        已超越美国模型（以 Qwen 超越 Llama 为标志）。
  \item DeepSeek-R1 是 2025 年全球被下载和部署次数最多的
        开放权重推理模型。
  \item 中国政府将"开源共享"写入了 2017 年《新一代人工智能发展规划》、
        "十四五"规划等顶层文件，形成了
        \textbf{自上而下的政策支持 + 自下而上的社区创新}的双轮驱动。
  \item 但风险也存在：开源模型的安全治理、出口管制合规性、
        以及"开放权重"（open-weight）与"真正开源"（open-source）
        之间的定义争议，仍是悬而未决的问题。
\end{itemize}
\end{corebox}


% ============================================================
%  §8.10  SYNTHESIS
% ============================================================

\section{总结：中国 AI 人才版图的深层逻辑}

纵观本章 130+ 位人物——从大模型到机器人，从自动驾驶到开源生态——
几个结构性特征清晰浮现：

\begin{corebox}[中国 AI 人才版图的七大特征]
\textbf{1. 学术-产业旋转门}。中国 AI 的几乎每一位重要人物都在
学术界和产业界之间经历过至少一次角色转换。唐杰（清华教授 $\to$ 智谱首席科学家）、
杨植麟（清华/CMU 学者 $\to$ 月之暗面 CEO）、印奇（清华姚班 $\to$ 旷视 CEO $\to$ 阶跃星辰董事长）、
何恺明（MSRA $\to$ FAIR $\to$ MIT 教授 + Google DeepMind）、
王鹤（北大教授 $\to$ 银河通用 CTO）——
这种高频率的角色转换是中国 AI 生态活力的根源。

\textbf{2. 清华-中科大双核}。在六小虎的创始人中，
唐杰/张鹏/王小川来自清华，闫俊杰来自中科大，
杨植麟是清华本科+CMU 博士。
在芯片领域，陈天石/陈云霁来自中科大少年班。
清华和中科大（尤其是少年班）构成了中国 AI 人才的双核心。

\textbf{3. 微软亚研院的"AI 黄埔军校"效应}。
姜大昕（阶跃星辰）、张正友（腾讯）、印奇（旷视）、
何恺明（ResNet）、曹旭东（Momenta）、袁进辉（硅基流动）——
微软亚洲研究院（MSRA）培养的 AI 人才遍布中国 AI 产业的每一个角落。
MSRA 自 1998 年成立以来，已成为中国 AI 人才储备的最大蓄水池之一。

\textbf{4. 代际更替加速}。从张钹院士（1935 年生）到
陈天石/闫俊杰（$\sim$1985--1990 年生），再到杨植麟（1992 年生）
和彭志辉（1993 年生），中国 AI 领军人物的平均年龄在快速下降。
DeepSeek 的核心团队平均年龄不到 30 岁。
这种代际更替的速度在全球 AI 产业中都属罕见。

\textbf{5. 出口管制的逆向激励}。美国芯片出口管制本意是限制中国 AI 能力，
但实际产生了复杂的逆向效应：(a) 催生了华为昇腾、寒武纪等替代方案的爆发；
(b) 迫使 DeepSeek 等公司在受限硬件上实现更高的工程效率；
(c) 推动了一批 NVIDIA/AMD 前高管回流中国创业；
(d) 硅基流动等公司率先实现国产芯片上的高效推理部署。
这是地缘政治时代技术竞争的典型"意外后果"。

\textbf{6. 从数字到物理的 AI 扩散}。中国 AI 不再局限于
大模型和聊天机器人——它正在向物理世界快速扩散。
宇树和智元的人形机器人上春晚、小马智行和文远知行登陆资本市场、
小鹏提出"Physical AI"战略——
2025--2026 年的关键词是\textbf{具身智能}和\textbf{AI 出行}。

\textbf{7. 开源作为国家战略}。与美国开源生态的"自下而上"特征不同，
中国的 AI 开源发展呈现"自上而下政策引导 + 自下而上社区创新"的双轮模式。
从 DeepSeek 到 Qwen，从 ChatGLM 到 InternLM，
中国开放权重模型在全球的分发量已超越美国——
这既是技术实力的体现，也是产业战略的结果。
\end{corebox}

\begin{warnbox}[中国 AI 人才面临的风险]
尽管中国 AI 人才储备在数量上已跻身世界前列，
但几个结构性风险值得关注：
\begin{itemize}[noitemsep]
  \item \textbf{顶尖人才外流}：何恺明、谢赛宁、陈丹琦等最具国际影响力的
        中国 AI 学者目前均任职于美国顶尖高校或实验室。
        中美地缘政治紧张可能进一步限制人才双向流动。
  \item \textbf{大厂人才虹吸}：六小虎和大厂 AI 实验室的高薪正在
        "掏空"高校的年轻研究者，可能影响中国 AI 的长期基础研究能力。
  \item \textbf{泡沫风险}：智谱 + MiniMax 合计市值超万亿港元，
        但合计营收可能不到 50 亿元（隐含市销率 200$\times$+）。
        估值泡沫一旦破裂，可能导致人才信心的连锁反应。
  \item \textbf{芯片持续受限}：尽管软件创新在一定程度上弥补了硬件劣势，
        但如果美国进一步收紧出口管制（如限制 HBM 供应或 EDA 工具），
        中国 AI 公司的训练能力仍可能面临瓶颈。
\end{itemize}
\end{warnbox}

这些人物的故事——从梁文锋的隐形帝国到汤晓鸥的未竟遗产，
从杨植麟的 AGI 梦想到陈天石兄弟的十年冷板凳，
从王兴兴 200 元造机器人到彭志辉从 B 站走向机器人工厂，
从楼天城的编程传说到曹旭东与全球车企的合作版图，
从袁进辉的推理引擎到邱锡鹏的开源先声——
共同编织了中国 AI 在 2020 年代的宏大叙事。

这个叙事的核心主题不是"追赶"，而是\textbf{在约束条件下的创造性突破}。
当算力受限时，用算法创新补偿；当生态封闭时，用开源打开局面；
当资本泡沫膨胀时，用工程效率回归本质；
当 AI 从数字世界溢出时，用制造业供应链拥抱物理世界。

这就是中国 AI 的东方战线——
从对话框到机器人关节，从推理 Token 到方向盘，
从 GitHub Star 到春晚舞台——
充满矛盾、充满张力，但也充满可能性。
